% =====================================================================
% Rapport de recherche — Master
% Pipeline « URDF vers robot commandé » par réseau de neurones informé
% par la physique (PINN gris) — Franka Emika Panda 7 DoF
% =====================================================================

\documentclass[a4paper, oneside]{article}

\usepackage[a4paper, fancysections, titlepage, pagenumber]{polytechnique}
\usepackage[french]{babel}
\usepackage[T1]{fontenc}
\usepackage[utf8]{inputenc}
\usepackage[hidelinks]{hyperref}
\usepackage{amsmath, amssymb}
\usepackage{booktabs}
\usepackage{array}
\usepackage{longtable}
% Unités physiques standardisées (Nm, Hz, rad, kg...) avec espacement et
% notation décimale françaises, cohérentes avec la convention à virgule
% utilisée par pgfplots ci-dessous.
\usepackage[locale=FR, per-mode=symbol]{siunitx}
\sisetup{output-decimal-marker={,}, group-separator={\,},
    list-final-separator={ et }, list-separator={, }}
\usepackage{tikz}
\usepackage{pgfplots}
\pgfplotsset{compat=1.16}
% Convention numérique française pour tous les graphiques
\pgfplotsset{
    /pgf/number format/use comma,
    /pgf/number format/1000 sep={\,},
}
\usetikzlibrary{arrows.meta, positioning, fit, backgrounds, calc}
\usepackage{caption}
\captionsetup{font=small, labelfont=bf}

% Bibliographie (BibTeX/Biber, remplace l'environnement thebibliography
% manuel). Style numérique pour préserver le rendu [1], [2]... déjà utilisé
% dans tout le texte par \cite{}. polytechnique.sty ne charge ni natbib ni
% biblatex et n'impose donc aucune convention de bibliographie particulière.
\usepackage{csquotes}
\usepackage[backend=biber, style=numeric, sorting=nyt, giveninits=true]{biblatex}
\addbibresource{bibliography.bib}

% Titre de paragraphe en vedette (run-in). Défini explicitement plutôt que
% d'utiliser \paragraph, car titlesec — chargé par polytechnique.sty — insère
% les \paragraph dans la table des matières en ignorant tocdepth. Une
% correction native à titlesec a été envisagée (reformater \paragraph via
% \titleformat/\titlespacing puis compter sur tocdepth pour filtrer l'entrée
% de table des matières), mais titlesec est connu pour ne pas toujours
% respecter tocdepth sur les niveaux non explicitement reformatés par
% \DeclareOption{fancysections} de polytechnique.sty ; sans vérification
% visuelle live de la table des matières sur les 40 occurrences de \parag du
% document, ce correctif plus « propre » n'a pas été appliqué pour ne pas
% risquer une régression silencieuse. Cette macro dédiée reste donc en place
% intentionnellement.
\newcommand{\parag}[1]{%
    \par\medskip\noindent\textbf{#1}\ \ignorespaces%
}

\newcommand{\code}[1]{%
    \mbox{\ttfamily\small%
        \detokenize{#1}%
    }%
}

% --- Notation (utilisée de façon cohérente dans tout le document) ---
\newcommand{\q}{\ensuremath{q}}
\newcommand{\qd}{\ensuremath{\dot{q}}}
\newcommand{\qdd}{\ensuremath{\ddot{q}}}
\newcommand{\taurbd}{\ensuremath{\tau_{\mathrm{rbd}}}}
\newcommand{\taumlp}{\ensuremath{\tau_{\mathrm{mlp}}}}
\newcommand{\taufric}{\ensuremath{\tau_{\mathrm{fric}}}}
\newcommand{\taures}{\ensuremath{\tau_{\mathrm{res}}}}
\newcommand{\taupred}{\ensuremath{\tau_{\mathrm{pred}}}}
\newcommand{\taureal}{\ensuremath{\tau_{\mathrm{real}}}}
\newcommand{\taucmd}{\ensuremath{\tau_{\mathrm{cmd}}}}
\newcommand{\Rset}{\ensuremath{\mathbb{R}}}

\author{Hugo Durieux}
\date{\today}
\title{Des fichiers URDF au robot commandé, en simulation}
\subtitle{Apprentissage de dynamique par réseau de neurones informé par la
physique et commande temps réel à \SI{1}{kHz} d'un Franka Emika Panda 7~DoF
simulé (Isaac~Sim / MuJoCo)}

\begin{document}

\maketitle

\begin{center}
\begin{minipage}{0.88\textwidth}
\small
\textbf{Résumé.} La commande précise d'un manipulateur repose sur un modèle
dynamique fidèle, mais le modèle rigide analytique issu d'un URDF est
systématiquement faux face aux frottements et à la compliance réels. Ce
travail construit et instrumente une chaîne entièrement automatisée, en
simulation, qui part du seul fichier URDF d'un Franka Emika Panda 7~DoF et
produit un robot simulé effectivement commandé par un modèle dynamique gris~:
un terme analytique calculé par RNEA (Pinocchio), jamais appris, composé à un
résidu neuronal contraint et conditionné par la charge portée, avec trois
degrés de garantie croissants (libre, imposé par la perte via un lagrangien
augmenté, structurel par construction pour la partie dissipative). Ce modèle
est inséré dans une loi de couple calculé plus PD à \SI{1000}{Hz} dont les gains
sont dérivés d'une borne d'erreur mesurée, pilotant un Franka simulé sous
MuJoCo via ROS2. Le modèle de référence atteint une erreur de prédiction de
couple de 0,30 à 0,90~Nm par axe (0,62 à 3,11\% de la limite de couple), avec
violation de limite de couple nulle et violation de dissipativité résiduelle
de $6\cdot 10^{-4}$ sur les données d'entraînement. Le suivi de trajectoire
par le modèle appris est confirmé numériquement en boucle fermée, mais la
tâche de préhension complète utilisée comme test d'intégration n'aboutit pas~:
ce rapport documente ce résultat négatif, l'immobilisation intermittente d'un
axe pendant l'exécution, dont une cause réelle mais partielle a été identifiée
et corrigée.

\medskip
\textbf{Abstract.} Accurate manipulator control requires a faithful dynamics
model, yet the analytical rigid-body model derived from a URDF is
systematically wrong once friction and transmission compliance are accounted
for. This work builds and instruments a fully automated, simulation-only
pipeline that takes as sole input the URDF file of a 7-DoF Franka Emika Panda
and produces a simulated robot actually driven by a learned grey-box dynamics
model~: an analytical term computed by RNEA (Pinocchio), never learned,
composed with a constrained, payload-conditioned neural residual exposing
three increasing degrees of guarantee (unconstrained, loss-imposed via an
augmented Lagrangian, structurally guaranteed by construction for the
dissipative part). This model is embedded in a \SI{1000}{Hz} computed-torque-plus-PD
control law whose gains are derived from a measured error bound, driving a
simulated Franka under MuJoCo through ROS2. The reference model reaches a
per-joint torque prediction error of 0.30 to 0.90~Nm (0.62 to 3.11\% of the
torque limit), with zero torque-limit violation and a residual dissipativity
violation of $6\cdot 10^{-4}$ on the training data. Closed-loop trajectory
tracking by the learned model is confirmed numerically, but the full grasping
task used as an integration test does not succeed~: this report documents
that negative result, an intermittent joint freeze during execution, for
which one real but partial cause was identified and fixed.
\end{minipage}
\end{center}

\setcounter{tocdepth}{2}
\tableofcontents
\newpage

% =====================================================================
\section{Introduction}
% =====================================================================

\subsection{Un modèle exact et néanmoins faux}

% TODO(hugo): confirmer le nom exact du laboratoire d'accueil et de l'équipe
% de recherche à Melbourne ; aucune trace de cette information n'a été
% trouvée dans le dépôt du projet (goal.md, PROJECT_STATE.md, PAPER_DRAFT.md)
% au moment de cette révision. Compléter la phrase ci-dessous en conséquence.
Ce travail a été mené dans le cadre d'un stage de recherche de Master, au sein
d'un laboratoire de robotique à Melbourne dont l'un des axes est l'interaction
humain-agent~: un robot qui partage un espace de travail avec une personne, ou
qui répond à une consigne de haut niveau plutôt qu'à une trajectoire
programmée, ne peut être sûr et réactif que si son modèle dynamique interne
est à la fois exact et interrogeable en temps réel. C'est cette exigence
double — précision du modèle, disponibilité en boucle de commande — qui motive
plus largement le choix d'un manipulateur physiquement instrumentable comme le
Franka Emika Panda comme cible de ce travail, au-delà du seul intérêt pour la
méthode d'apprentissage elle-même.

La commande précise d'un bras manipulateur repose sur la connaissance de sa
dynamique, c'est-à-dire de l'application qui relie une trajectoire articulaire
souhaitée aux couples moteurs qui la réalisent. Pour un robot série rigide à
$n$ degrés de liberté, cette relation prend la forme classique du modèle
dynamique inverse
%
\begin{equation}
    \tau = M(\q)\,\qdd + C(\q,\qd)\,\qd + G(\q) ,
    \label{eq:rbd}
\end{equation}
%
où $\q, \qd, \qdd \in \Rset^{n}$ désignent respectivement les positions,
vitesses et accélérations articulaires, $M(\q)$ la matrice d'inertie,
$C(\q,\qd)$ la matrice des termes de Coriolis et centrifuges, $G(\q)$ le
vecteur de gravité et $\tau \in \Rset^{n}$ le vecteur des couples
articulaires. L'équation~\eqref{eq:rbd} est entièrement déterminée dès lors que
la géométrie et les paramètres inertiels du robot sont connus; elle se calcule
efficacement par l'algorithme récursif de Newton-Euler (\emph{Recursive
Newton-Euler Algorithm}, RNEA) à partir d'une description cinématique et
inertielle, typiquement un fichier URDF (\emph{Unified Robot Description
Format}).

Le problème est que l'équation~\eqref{eq:rbd} est fausse sur un robot réel. Les
frottements secs et visqueux des réducteurs harmoniques, la compliance des
transmissions, le jeu mécanique (\emph{backlash}), les pertes dépendantes de la
température et l'usure des actionneurs produisent un écart systématique entre
le couple prédit par le modèle rigide et le couple réellement mesuré. Cet
écart, appelé ici couple résiduel, est de l'ordre de plusieurs newton-mètres
sur un Franka Emika Panda et n'est pas modélisable analytiquement de manière
fiable. C'est précisément là que l'apprentissage automatique intervient — et
c'est aussi là que se joue tout le contenu scientifique du présent travail:
non pas dans le remplacement de~\eqref{eq:rbd}, mais dans le traitement de ce
qu'elle ne dit pas.

Deux voies s'opposent. La première apprend la dynamique de bout en bout, comme
une boîte noire: un réseau de neurones reçoit l'état et prédit le couple.
Cette approche est flexible mais coûteuse en données, ne garantit aucune
propriété physique et généralise mal hors de la distribution d'entraînement. La
seconde, dite \emph{informée par la physique}, contraint le réseau à respecter
la structure des équations de la mécanique — soit en inscrivant cette structure
dans l'architecture, soit en l'imposant par des termes de perte. Les réseaux
lagrangiens profonds~\cite{lutter2019iclr} et, plus récemment, les travaux de
Liu, Borja et Della~Santina~\cite{liu2024} ont établi la viabilité théorique et
expérimentale de cette seconde voie sur des manipulateurs réels.

\subsection{Objectif et verrou scientifique}
\label{sec:verrou}

L'objectif de ce travail est de \textbf{construire et valider une chaîne
automatique qui prend en entrée le seul fichier de description d'un
manipulateur et produit en sortie un robot effectivement commandé} par un
modèle dynamique appris, interrogé à chaque période de la boucle de commande.
Cette chaîne comprend l'apprentissage d'un modèle dynamique gris
(\emph{grey-box}), sa mise en service dans une loi de commande temps réel, et
une tâche de préhension complète servant de test d'intégration de bout en bout.

Formulé ainsi, l'objectif paraît relever de l'ingénierie. Trois difficultés
distinctes le rendent non trivial, et ce sont elles qui définissent la
contribution.

\parag{La généralité s'obtient contre la précision.}
La littérature obtient ses meilleurs résultats en dérivant à la main les
équations d'un robot particulier, puis en ajustant un réseau sur cette
structure. Rendre le procédé automatique à partir d'un simple URDF signifie
renoncer à toute dérivation manuelle: le graphe cinématique, les dimensions du
réseau, l'excitation des trajectoires d'apprentissage et les bornes physiques
doivent être extraits du fichier de description. Or ce fichier est lui-même une
source d'erreurs silencieuses. Au cours de ce travail, l'URDF de référence du
projet s'est révélé dépourvu de toute balise inertielle: RNEA calculait depuis
des inerties de remplacement, et non depuis les propriétés massiques réelles du
Franka, pendant toute la première partie du projet
(section~\ref{sec:disc-urdf}).

\parag{Les garanties physiques survivent mal à l'apprentissage.}
Un résidu appris librement peut injecter de l'énergie dans le système —
physiquement absurde pour un terme censé représenter des frottements — et
prédire des couples supérieurs aux limites des actionneurs. Sur un système à
7~degrés de liberté, ces contraintes doivent être satisfaites simultanément sur
sept axes dont les limites de couple diffèrent d'un facteur~7 (\SI{87}{Nm}
pour les quatre premiers axes, \SI{12}{Nm} pour les trois derniers), ce qui
déséquilibre
naturellement l'optimisation. Les preuves de concept de la littérature sont le
plus souvent validées sur des pendules ou des bras plans à deux axes; le
passage à 7~DoF n'est pas une simple mise à l'échelle.

\parag{L'exactitude hors ligne ne fait pas une commande.}
Une erreur de prédiction faible sur un jeu de validation ne dit rien de la
stabilité en boucle fermée. Insérer un réseau de neurones dans une boucle à
\SI{1000}{Hz} impose une inférence sous la milliseconde, un ordonnancement temps réel
fiable, et une loi de commande dont on sache borner l'erreur en présence
d'erreur de modèle. C'est la difficulté qui a produit le plus de défauts réels:
sur l'ensemble des sessions de mise au point documentées, la quasi-totalité des
bogues bloquants se situait non pas dans les mathématiques de la commande, mais
à l'interface entre le modèle appris, l'intergiciel ROS2 et le simulateur
(sections~\ref{sec:res-loop} et~\ref{sec:res-task}).

\subsection{Contributions et état d'avancement}
\label{sec:contributions}

Ce travail se positionne explicitement par rapport à l'état de l'art défini par
Liu \emph{et al.}~\cite{liu2024}, référence la plus proche: même robot (Franka
Panda 7~DoF), même classe de méthodes (réseaux informés par la physique), même
finalité (modéliser puis commander). Quatre contributions sont revendiquées.

\begin{enumerate}
    \item \textbf{Une chaîne automatisée \og URDF vers modèle \fg{}.} Le
    fichier URDF est la seule entrée: le graphe cinématique et les paramètres
    inertiels en sont extraits par Pinocchio, le terme analytique est évalué par
    RNEA sans aucune dérivation manuelle, l'architecture du résidu est
    dimensionnée depuis le nombre d'axes, et les bornes physiques (couples,
    vitesses) alimentent directement les contraintes d'entraînement. Liu
    \emph{et al.} dérivent au contraire les équations de leur réseau
    d'énergie spécifiquement pour leur plateforme.

    \item \textbf{L'intégration du modèle appris dans une boucle de commande à
    \SI{1000}{Hz}.} Le modèle n'est pas évalué seulement hors ligne sur des données
    enregistrées: il est interrogé à chaque période de commande par un n\oe ud
    ROS2 qui pilote un Franka simulé sous MuJoCo à travers
    \code{ros2_control}. Liu \emph{et al.} s'arrêtent à \SI{500}{Hz}.

    \item \textbf{La montée en dimension à 7~DoF sous contraintes physiques
    dures.} Les contraintes de limites de couple et de dissipativité sont
    imposées simultanément sur les sept axes par un lagrangien augmenté à
    montée duale, et complétées par une garantie structurelle de dissipativité
    obtenue par construction (matrice de frottement diagonale définie
    positive), et non par pénalisation.

    \item \textbf{Un protocole, implémenté mais non encore validé, de
    réduction de l'écart simulation-réalité.} Le modèle est pré-entraîné
    massivement en simulation, puis affiné sur un petit enregistrement réel
    avec toutes les couches gelées sauf les deux dernières, sous la même perte
    informée par la physique. Ce protocole est correctement implémenté et
    vérifié sur le plan physique (section~\ref{sec:finetune}), mais il n'a
    jamais été exécuté sur un enregistrement réel~: à ce titre, il doit être
    lu comme une perspective outillée plutôt que comme un résultat au même
    titre que les contributions~1 à~3, distinction reprise explicitement dans
    le paragraphe qui suit et dans les limites (section~\ref{sec:limites}).
\end{enumerate}

Il faut annoncer d'emblée l'état réel d'avancement: les contributions~1 et~3
sont implémentées et validées quantitativement; la contribution~2 est validée
en boucle fermée mais la mesure de latence d'inférence qui la substantierait
pleinement n'a pas été réalisée. La quatrième reste, par construction, d'une
nature différente des trois premières~: le protocole de transfert
simulation-réalité est implémenté et vérifié sur le plan physique, mais aucun
enregistrement sur robot réel n'a été effectué, si bien qu'il n'a en réalité
jamais été mis à l'épreuve. Un protocole non testé ne constitue pas un
résultat validé~; il est présenté ici comme une contribution \emph{d'outillage
et de conception}, dont la validation reste entièrement à faire, et non comme
une contribution empiriquement établie au même titre que les trois autres. Ce
rapport rend compte du travail effectué, y compris de ce qui n'a pas
fonctionné: la tâche de préhension, qui devait servir de démonstration
applicative, n'aboutit pas, et l'analyse de son échec constitue le principal
résultat négatif du travail.

La suite est organisée de manière à faire porter chaque section une part de
l'argument. La section~\ref{sec:etat-art} construit, à partir des travaux
existants, le verrou précis que ce travail attaque, et le formalise sous forme
d'un tableau d'écarts avec la référence. La section~\ref{sec:methode} déduit de
ce verrou une méthode et justifie chacun de ses choix, y compris ceux qui
consistent à rejeter un mécanisme attractif de la littérature. La
section~\ref{sec:resultats} confronte chaque contribution revendiquée à une
mesure susceptible de la démentir. La section~\ref{sec:discussion} interprète
ces mesures, y compris les défavorables, et énonce les limites. La
section~\ref{sec:conclusion} conclut.

% =====================================================================
\section{Où placer la connaissance physique~? Lecture critique des
travaux existants}
\label{sec:etat-art}
% =====================================================================

Une trentaine d'articles ont été dépouillés de manière systématique pendant ce
travail. Chacun a donné lieu à une extraction structurée selon une grille
constante — résumé, plateforme, résultats chiffrés, nouveautés candidates,
puis verdict motivé de reprise, de rejet ou de mise en attente — consignée dans
un registre unique et associée, pour les nouveautés retenues, à une trace
d'implémentation. Cette discipline importe pour la lecture qui suit: plusieurs
mécanismes discutés ci-dessous ont été rejetés, et le rejet est alors argumenté
par un critère explicite plutôt que passé sous silence.

Reste à choisir un axe de lecture. Regrouper ces articles par thème applicatif
serait sans portée, car ils visent presque tous la même chose. L'axe retenu ici
est celui qui sépare réellement les méthodes entre elles: \textbf{où la
connaissance physique est injectée dans le problème d'apprentissage, et ce
qu'elle garantit en retour}. Un modèle appris comporte en effet trois lieux
d'insertion possibles — ses paramètres, son architecture, et sa fonction
objectif — et le choix du lieu détermine la nature de la garantie obtenue:
paramètres physiquement admissibles, propriété vraie par construction, ou
contrainte satisfaite en moyenne sur les données. Cette lecture par lieu
d'insertion est celle que retient également la revue de
Ma \emph{et al.}~\cite{ma2025}, qui distingue entrées informées, pertes
informées et architectures informées, en la justifiant par l'anatomie même d'un
modèle d'apprentissage. Les sections~\ref{sec:sota-param}
à~\ref{sec:sota-compo} parcourent cet axe. La
section~\ref{sec:sota-boucle} introduit une contrainte d'un autre ordre, que
l'axe précédent ignore: ce que la mise en boucle fermée exige en retour du
modèle. La section~\ref{sec:sota-gap} en déduit le verrou.

\subsection{La connaissance dans les paramètres: identifier un modèle rigide}
\label{sec:sota-param}

Le premier lieu d'insertion est le plus ancien: conserver intégralement la
structure analytique~\eqref{eq:rbd} et ne chercher qu'à en identifier les
paramètres. On excite le robot par des trajectoires riches, puis on résout un
problème de moindres carrés sur les paramètres inertiels.
Sutanto \emph{et al.}~\cite{sutanto2020} en proposent la version moderne et
différentiable: RNEA est réécrit comme un graphe de calcul PyTorch, et les
paramètres inertiels de chaque corps deviennent des variables optimisables.
Leur apport propre est instructif pour le présent travail, car il porte
précisément sur la \emph{forme} de la garantie: ils comparent une série de
re-paramétrisations de plus en plus contraintes du tenseur d'inertie — non
structurée, symétrique, définie positive, respectant l'inégalité triangulaire
des moments principaux — et montrent que les variantes les plus contraintes
convergent en 2~époques contre 61 pour la variante libre, tout en généralisant
mieux à des tâches non vues. Une contrainte plus forte ne coûte donc pas
nécessairement en performance: elle peut au contraire réduire le coût
d'apprentissage.

L'intérêt de cette approche est sa parcimonie et son interprétabilité: le
modèle obtenu est celui de la mécanique, avec des paramètres qui ont un sens
physique. Sa limite est structurelle, et non affaire de réglage: quelle que
soit la qualité de l'identification, un modèle rigide ne peut représenter ni
frottements non linéaires ni compliance, car ces phénomènes ne figurent pas
dans son espace de fonctions. Sutanto \emph{et al.} le constatent eux-mêmes en
passant du simulateur au robot réel, où le frottement statique reste le défaut
résiduel de leur formulation.

\subsection{La connaissance dans l'architecture: apprendre une structure
énergétique}
\label{sec:sota-archi}

Le deuxième lieu d'insertion déplace la physique dans l'architecture du réseau
lui-même. Les réseaux lagrangiens profonds (DeLaN) de Lutter, Ritter et
Peters~\cite{lutter2019iclr} apprennent la matrice d'inertie $M(\q)$ sous forme
factorisée de Cholesky $M = LL^{\top}$, avec une diagonale passée par une
fonction Softplus, ce qui garantit la définie-positivité en tout point de
l'espace articulaire; la gravité est apprise par une tête séparée, et
l'équation d'Euler-Lagrange est imposée par construction. Le même groupe étend
ensuite l'approche à la commande fondée sur l'énergie pour des systèmes
sous-actionnés~\cite{lutter2019iros}. Duong \emph{et al.}~\cite{duong2024}
transposent l'idée dans le formalisme port-hamiltonien sur groupes de Lie, avec
des sous-réseaux distincts et factorisés par Cholesky pour l'inertie, la
dissipation, l'énergie potentielle et la matrice d'entrée.
Wang, Chen et Ding~\cite{wang2025spel} poursuivent avec une formulation
hamiltonienne symplectique et apportent un argument géométrique dont ce travail
fera un usage direct: pour des liaisons rotoïdes tournant chacune autour d'un
axe unique, la matrice de dissipation par corps est de rang~1, ce qui autorise
une structure creuse.

Liu, Borja et Della~Santina~\cite{liu2024} constituent la référence directe de
ce travail. Ils étendent les réseaux lagrangiens et hamiltoniens à la
dissipation et à l'actionnement non colocalisé, et — c'est leur apport le plus
utile ici — démontrent la stabilité au sens de Lyapunov d'un correcteur à
couple calculé augmenté d'un terme proportionnel-dérivé lorsque l'erreur de
modèle est bornée. Leur validation est expérimentale sur un Franka Panda réel:
erreur de prédiction ramenée de 182{,}5 à 2{,}68~rad$^2$ face à une référence
boîte noire, avec seulement 25\,000 échantillons d'apprentissage contre
550\,000 pour cette référence, et commande effective du robot réel à \SI{500}{Hz}.

La force de ce lieu d'insertion est que la garantie ne dépend pas des données:
une inertie factorisée par Cholesky est définie positive quel que soit le jeu de
poids. Sa faiblesse est le prix payé pour cette force. Apprendre $M(\q)$
consiste à dépenser de la capacité de réseau — et des données — pour
reconstruire une quantité que l'URDF fournit déjà exactement, et à perdre au
passage la propriété qui intéresse ce travail: qu'un modèle valide s'obtienne
automatiquement pour n'importe quelle description de robot correctement
renseignée. Ces travaux paient donc en généralité automatique ce qu'ils gagnent
en garantie.

Ce constat mérite d'être rendu explicite pour le formalisme port-hamiltonien de
Duong \emph{et al.}~\cite{duong2024} en particulier, plutôt que laissé comme un
cas parmi d'autres de ce lieu d'insertion. La passivité y est garantie de
façon intrinsèque et inconditionnelle par la structure hamiltonienne elle-même
— une propriété strictement plus forte que la contrainte de dissipativité
imposée ici par pénalisation duale au degré~1 (section~\ref{sec:degre1}), qui
n'est vraie qu'en moyenne sur les données vues à l'entraînement. Ce formalisme
a été examiné et écarté comme pilier architectural pour deux raisons propres à
ce travail, distinctes du seul argument de généralité automatique ci-dessus.
D'abord, une architecture hamiltonienne sur groupes de Lie, comme celle de
Duong \emph{et al.}, doit réapprendre l'inertie, la dissipation, l'énergie
potentielle et la matrice d'entrée depuis les données~: exactement ce que le
verrou de la section~\ref{sec:sota-gap} interdit, puisque ces quantités sont
déjà disponibles exactement via RNEA pour tout URDF correctement renseigné.
Ensuite, et c'est un argument de conception plus ciblé~: la propriété qui
intéresse spécifiquement ce travail n'est pas la passivité de l'ensemble du
système appris, mais la seule non-injection d'énergie par le terme dissipatif
du résidu — une propriété strictement plus étroite, pour laquelle une
architecture hamiltonienne complète est disproportionnée. Le degré~2 de la
section~\ref{sec:frictionnet} obtient précisément cette garantie plus étroite,
par construction et de façon inconditionnelle, avec une matrice de frottement
diagonale à un coût architectural bien moindre qu'un réseau hamiltonien
complet. Le prix de ce choix est net et assumé~: contrairement au formalisme
port-hamiltonien, rien ne garantit ici la passivité du couple total
$\taupred$, seul le terme $\taufric$ l'est structurellement — c'est pourquoi la
contrainte de dissipativité~\eqref{eq:gdissip} du degré~1 reste nécessaire pour
couvrir $\taumlp$, dont la sortie demeure libre.

Un second écart, plus discret, mérite d'être relevé ici car il structure toute
la méthode de la section~\ref{sec:methode}. Liu \emph{et al.} évitent
délibérément toute mesure d'accélération, car $\qdd$ est bruité sur robot réel;
ils intègrent donc le modèle dans le temps, ce qui les conduit à une perte de
\emph{rollout} par intégration Runge-Kutta d'ordre~4. Ce choix est cohérent chez
eux et incompatible avec l'architecture retenue ici, où $\qdd$ n'intervient que
hors ligne — sous forme analytique lors de la préparation du jeu de données — ou
sous sa forme \emph{désirée}, issue du planificateur, dans la boucle de
commande. À aucun moment une accélération mesurée n'est employée.

\subsection{La connaissance partagée: composer un terme analytique et un
résidu contraint}
\label{sec:sota-compo}

Le troisième lieu d'insertion n'est ni les paramètres seuls ni l'architecture
seule: il compose un terme analytique connu avec un résidu appris, et impose
les propriétés physiques restantes par la fonction objectif. Djeumou
\emph{et al.}~\cite{djeumou2022} en donnent la formulation de référence: des
termes de champ de vecteurs connus sont combinés à des résidus neuronaux, et
les contraintes physiques sont imposées par un lagrangien augmenté avec montée
duale sur les multiplicateurs. Leurs résultats séparent nettement les deux
effets: gain d'un facteur~100 sur l'erreur de rollout face à une boîte noire
grâce à la composition, et réduction de la violation de contraintes de quatre
ordres de grandeur grâce au lagrangien augmenté — c'est-à-dire que la
composition apporte la précision et que le mécanisme dual apporte le respect de
la physique, les deux n'étant pas interchangeables.

Deng \emph{et al.}~\cite{deng2024} suivent une logique voisine mais plus
prescriptive: les sous-termes de la dynamique inverse (produits
trigonométriques, produits polynomiaux, combinaisons linéaires) sont encodés
comme des couches structurelles dures, ce qui rend les poids du réseau
interprétables comme paramètres physiques et permet de les comparer aux valeurs
vraies. Leur validation, sur le dernier axe d'un KUKA iiwa 7~DoF, contient un
avertissement utile: le modèle de référence purement appris atteint une erreur
comparable à la leur \emph{en moyenne} tout en produisant une erreur maximale de
44{,}4~Nm, c'est-à-dire des prédictions physiquement absurdes que la moyenne
masque. Li \emph{et al.}~\cite{li2025compliant} appliquent une architecture
grise voisine à des manipulateurs compliants et confirment indépendamment la
supériorité du gris sur le noir dans un autre domaine. La revue de
Toscano \emph{et al.}~\cite{toscano2025} recense huit années d'avancées et
établit que les compositions grises, l'imposition de contraintes par lagrangien
augmenté et l'affinage à ossature gelée constituent des pratiques établies
plutôt que des paris.

Deux résultats de cette lecture bornent par avance ce que le présent travail
peut prétendre. Le premier est l'avertissement de Deng \emph{et al.} sur les
métriques moyennes, qui impose de rapporter des queues de distribution et non
seulement des erreurs quadratiques moyennes; il est suivi en
section~\ref{sec:res-bound}. Le second est le résultat négatif de
Prabhakar \emph{et al.}~\cite{prabhakar2026}: sur des dynamiques de contact,
les contraintes physiques \emph{dégradent} l'extrapolation temporelle au-delà de
la fenêtre d'apprentissage, de 6 à 10\% d'erreur supplémentaire. Ce résultat
n'est pas contourné dans ce travail mais assumé comme une restriction de portée
déclarée (section~\ref{sec:limites}).

C'est de cette manière de partager la connaissance que le présent travail est
le plus proche: l'architecture développée en section~\ref{sec:methode} en est
une instance. Il convient de le dire sans ambiguïté, car deux nouveautés
candidates issues de Djeumou \emph{et al.} — la structure compositionnelle et le
lagrangien augmenté à montée duale — ont été examinées puis classées \og déjà
présentes \fg{}: toutes deux étaient implémentées indépendamment avant la
lecture de l'article, qui joue donc le rôle d'une validation externe de
l'architecture et non d'une source. La contribution propre de ce travail est
ailleurs, et se révèle en changeant de question.

\subsection{Ce que la fermeture de la boucle impose en retour}
\label{sec:sota-boucle}

Les trois sections précédentes examinent comment la physique entre dans le
modèle. Une contrainte d'un tout autre ordre apparaît lorsqu'on demande au
modèle de piloter un robot: la boucle fermée impose ses propres exigences en
retour, et elles éliminent des mécanismes qui étaient pourtant les meilleurs sur
le plan de la modélisation. C'est un enseignement central de la revue
d'Ibarz \emph{et al.}~\cite{ibarz2021}, qui consacre à la latence et à
l'exécution asynchrone un traitement de premier plan: un décalage de quelques
millisecondes entre mesure et actionnement suffit à faire échouer des
algorithmes par ailleurs à l'état de l'art, parce qu'il rompt l'hypothèse
d'exécution synchrone sur laquelle ils reposent.

\parag{La cadence disqualifie des mécanismes.} Les intégrations temps réel
publiées se situent toutes entre 2 et \SI{100}{Hz}. Wang \emph{et al.}~\cite{wang2024cac}
embarquent un PINN discret dans une commande prédictive non linéaire pour un
bras multi-axes, avec un résidu physique intégré par la méthode des trapèzes à
\SI{20}{Hz} et un filtre de Kalman étendu à l'interface. Li \emph{et al.}~\cite{li2025ral}
opèrent à \SI{100}{Hz}. Djeumou \emph{et al.}~\cite{djeumou2024drift} proposent, pour
la conduite à la limite d'adhérence, une architecture à cadences découplées: un
modèle de diffusion lent à \SI{2}{Hz} estime les paramètres de frottement, tandis
qu'une commande prédictive rapide s'exécute par-dessus. Viser \SI{1000}{Hz}, comme
ici, rend inutilisables le résidu discrétisé à \SI{20}{Hz} et le filtre de Kalman
étendu, dont l'inversion matricielle par pas de commande est incompatible avec
un budget de l'ordre de la milliseconde. La cadence n'est donc pas un paramètre
de réglage: c'est une contrainte qui sélectionne l'architecture.

\parag{La charge portée doit être traitée quelque part.} Une masse portée
modifie l'inertie, la gravité et les frottements perçus. Li \emph{et al.}~\cite{li2025ral}
l'identifient en ligne à partir des courants moteurs d'actionneurs à
entraînement quasi direct, puis injectent un corps virtuel correspondant dans
le modèle Pinocchio, lequel alimente une commande prédictive non linéaire.
Hu \emph{et al.}~\cite{hu2026} traitent par une approche PINN l'erreur induite
par le jeu mécanique sur des robots industriels et établissent un fait qui
oriente directement la décomposition retenue ici: l'erreur de \emph{backlash}
dépend de la direction de la vitesse et de la position articulaire, mais non de
la masse portée. Il existe donc une frontière naturelle entre ce qui dépend de
la charge et ce qui dépend de l'état de mouvement. Aucun des travaux cités ne
l'exploite dans les deux sens à la fois.

\parag{Le transfert et la planification sont des dépendances, non des
contributions.} Sur le versant planification, Ni et Qureshi~\cite{ni2024} et
Liu \emph{et al.}~\cite{liu2024ntfields} utilisent des PINN résolvant
l'équation d'Eikonal pour planifier sans données d'expert, avec des temps de
planification de 0{,}05 à 0{,}15~s contre 0{,}06 à 6{,}77~s pour des références
classiques. Sur le versant transfert, Ibarz \emph{et al.}~\cite{ibarz2021}
documentent le coût du passage simulation-réalité,
Duong \emph{et al.}~\cite{duong2024} proposent un protocole d'affinage court à
ossature gelée, et Chen \emph{et al.}~\cite{chen2025soft} synthétisent les
approches pilotées par les données pour les robots mous, où le même problème se
pose de manière aiguë. Enfin, Fang \emph{et al.}~\cite{fang2026} et
Agrawal \emph{et al.}~\cite{agrawal2026} illustrent l'automatisation croissante
de la chaîne elle-même, respectivement pour la cinématique inverse de
manipulateurs mobiles redondants et pour l'orchestration automatisée de
l'entraînement puis du déploiement d'un PINN sous ROS2. Ce dernier travail est
instructif par sa manière de rendre compte du coût de déploiement: il publie
une décomposition de la latence de bout en bout composant par composant, où
l'inférence du réseau ne pèse que 2~ms sur 40{,}5~ms au total, le reste étant
absorbé par la perception et le passage d'interface. Une chaîne complète se juge
donc sur son maillon lent, pas sur son réseau.

\subsection{Le verrou retenu, et le positionnement de ce travail}
\label{sec:sota-gap}

La lecture qui précède permet d'énoncer le verrou précisément. Les approches qui
placent la physique dans les paramètres obtiennent un modèle interprétable mais
structurellement incapable de représenter les frottements. Celles qui la
placent dans l'architecture obtiennent des garanties inconditionnelles, mais
paient en réapprenant ce que la description du robot fournit déjà, ce qui
interdit l'automatisation depuis un simple URDF. Celles qui la partagent entre
un terme analytique et un résidu contraint évitent ces deux écueils, mais
aucune, à notre connaissance, ne conditionne le résidu par la charge portée ni
ne ferme la boucle sur un manipulateur industriel à sept axes à une cadence
compatible avec un asservissement de couple. Autrement dit: \textbf{la
composition grise est un acquis, mais elle n'a pas encore été portée
simultanément à sept axes, sous contraintes physiques dures, conditionnée par la
charge, et interrogée à chaque période d'une boucle à \SI{1}{kHz}.} C'est cet
ensemble de conditions simultanées, et non l'un de ses éléments pris isolément,
qui constitue l'objet de ce travail.

Le tableau~\ref{tab:gaps} exprime ce verrou sous la forme d'écarts point par
point avec la référence la plus proche~\cite{liu2024}. Il fournit la définition
opérationnelle des quatre contributions annoncées en
section~\ref{sec:contributions}, et c'est à lui que la
section~\ref{sec:resultats} confrontera les mesures.

\begin{table}[htbp]
\centering
\caption{Positionnement par rapport à la référence Liu \emph{et al.}
(2024)~\cite{liu2024}, article le plus proche de ce travail (même robot, même
classe de méthodes, même finalité).}
\label{tab:gaps}
\small
\begin{tabular}{@{}p{0.44\textwidth}p{0.48\textwidth}@{}}
\toprule
\textbf{Limite de Liu \emph{et al.} (2024)} & \textbf{Réponse de ce travail} \\
\midrule
Équations dérivées manuellement pour la plateforme visée & Chaîne automatisée
URDF $\rightarrow$ Pinocchio $\rightarrow$ réseau, sans dérivation manuelle
(contribution~1) \\
Commande à \SI{500}{Hz}, sans intégration dans un solveur & Boucle ROS2 à \SI{1000}{Hz}
interrogeant le modèle appris à chaque pas (contribution~2) \\
Aucun conditionnement par la charge portée & Charge $\delta$ traitée par deux
voies: inertie de l'effecteur dans RNEA, et entrée explicite du résidu;
testée à \SIlist{0;1;3}{kg} \\
Aucun protocole de transfert simulation-réalité & Pré-entraînement en
simulation puis affinage à ossature gelée sous perte informée
(contribution~4) \\
Aucune limite de couple ni de vitesse dans la perte & Lagrangien augmenté sur
les limites de couple et la dissipativité, plus écrêtage dur à la publication
des consignes \\
Réseau d'énergie intégralement appris, sans terme RNEA analytique &
Architecture grise: RNEA jamais modifié, seul le résidu est appris
(contribution~3) \\
\bottomrule
\end{tabular}
\end{table}

% =====================================================================
\section{Méthode}
\label{sec:methode}
% =====================================================================

\subsection{Décomposition: ce qui est calculé, ce qui est appris}
\label{sec:overview}

Toute la méthode découle d'une seule décision, prise en réponse au verrou de la
section~\ref{sec:sota-gap}: la dynamique du corps rigide est \emph{calculée},
jamais apprise, et l'apprentissage est confiné à ce qu'elle ne couvre pas. Le
modèle prédit donc
%
\begin{equation}
    \taupred = \underbrace{\taurbd(\q,\qd,\qdd,\delta)}_{\text{calculé
    (RNEA, URDF)}}
    \;+\;
    \underbrace{\taures(\q,\qd,\delta)}_{\text{appris, contraint}} ,
    \label{eq:decompo}
\end{equation}
%
où $\delta \in \Rset_{+}$ désigne la masse portée par l'effecteur. Cette
frontière n'est pas arbitraire: elle suit celle qu'établit
Hu \emph{et al.}~\cite{hu2026}, à savoir que les effets non modélisés dépendent
de l'état de mouvement mais non de la charge, tandis que les effets de charge
sont, eux, exactement ceux que le modèle rigide sait traiter. Chaque terme
de~\eqref{eq:decompo} reçoit donc la part du problème pour laquelle il est
qualifié, sans recouvrement.

La figure~\ref{fig:pipeline} situe les deux régimes d'utilisation de ce modèle.
Hors ligne, un unique fichier URDF alimente deux consommateurs: Pinocchio, qui
en construit le modèle rigide et évalue RNEA, et le simulateur de génération de
données, qui exécute des trajectoires d'excitation et mesure les couples
effectifs; la différence entre couple mesuré et couple calculé constitue la
cible d'apprentissage du résidu. En ligne, le modèle entraîné est chargé par un
n\oe ud ROS2 qui l'interroge à chaque pas de commande, dans une loi de couple
calculé pilotant un Franka simulé.

\begin{figure}[htbp]
\centering
\begin{tikzpicture}[
    font=\small,
    box/.style={draw, rounded corners=2pt, minimum height=8mm, align=center,
                inner sep=3pt, fill=white},
    white/.style={box, fill=black!4},
    learn/.style={box, fill=black!12},
    rt/.style={box, fill=black!7},
    ar/.style={-{Latex[length=2mm]}, thick},
    node distance=4mm
]
% --- entrée
\node[box] (urdf) {URDF\\\code{panda.urdf}};

% --- régime hors ligne
\node[white, right=9mm of urdf, yshift=9mm] (pin) {Pinocchio\\RNEA \eqref{eq:rbd}};
\node[white, right=9mm of urdf, yshift=-9mm] (sim) {excitation Sobol\\+ Fourier
(Isaac~Sim)};
\node[box, right=9mm of pin, yshift=-9mm, minimum height=10mm] (res)
    {cible du résidu\\$\taureal-\taurbd$};
\node[learn, right=9mm of res] (train) {entraînement\\MSE + lagrangien
augmenté};
\node[box, right=9mm of train] (ckpt) {modèle\\$\theta^{*}$};

% --- régime en ligne
\node[rt, below=14mm of train, xshift=-16mm] (moveit) {planification\\MoveIt2};
\node[rt, right=8mm of moveit] (ctrl) {couple calculé + PD\\\SI{1000}{Hz}};
\node[rt, right=8mm of ctrl] (robot) {Franka simulé\\MuJoCo};

\draw[ar] (urdf.east) -- (pin.west);
\draw[ar] (urdf.east) -- (sim.west);
\draw[ar] (pin.east) -- (res.west);
\draw[ar] (sim.east) -- (res.west);
\draw[ar] (res) -- (train);
\draw[ar] (train) -- (ckpt);
\draw[ar] (ckpt.south) -- ++(0,-7mm) -| (ctrl.north);
\draw[ar] (moveit) -- (ctrl);
\draw[ar] (ctrl) -- (robot);
\draw[ar] (robot.south) -- ++(0,-5mm) -| node[below, pos=0.25]
    {$\q,\qd$ mesurés} (ctrl.south);
\draw[ar] (urdf.south) |- ++(0,-30mm) -| (moveit.south);
\end{tikzpicture}
\caption{Les deux régimes d'utilisation du modèle. Le fichier URDF est la seule
entrée: il alimente le calcul analytique (Pinocchio/RNEA), la génération de
données d'excitation, et la planification. Rangée du haut: identification hors
ligne du résidu. Rangée du bas: boucle de commande temps réel, qui interroge le
modèle entraîné à chaque période. Les blocs clairs sont analytiques, le bloc
foncé est appris.}
\label{fig:pipeline}
\end{figure}

Le robot cible est un Franka Emika Panda à 7~degrés de liberté. Ses constantes
physiques, utilisées de manière centralisée par tous les modules, sont
rappelées au tableau~\ref{tab:constantes}. L'asymétrie d'un facteur~7 entre les
limites de couple des quatre premiers axes et des trois derniers y mérite une
attention particulière: elle intervient trois fois dans la suite, comme borne
de contrainte d'entraînement (section~\ref{sec:degre1}), comme source du
déséquilibre d'optimisation discuté en section~\ref{sec:res-desequilibre}, et
comme dénominateur de la seule grandeur d'erreur réellement interprétable pour
la commande (section~\ref{sec:res-quoi}).

\begin{table}[htbp]
\centering
\caption{Constantes physiques du Franka Emika Panda utilisées dans tout le
travail (source: fiche technique constructeur, centralisée dans
\code{network/constants.py}).}
\label{tab:constantes}
\small
\begin{tabular}{@{}lccccccc@{}}
\toprule
Axe $j$ & 1 & 2 & 3 & 4 & 5 & 6 & 7 \\
\midrule
Limite de couple $\bar{\tau}_j$ [\si{Nm}] & \num{87} & \num{87} & \num{87} &
\num{87} & \num{12} & \num{12} & \num{12} \\
Vitesse maximale [\si{\radian\per\second}] & \num{2,175} & \num{2,175} &
\num{2,175} & \num{2,175} & \num{2,610} & \num{2,610} & \num{2,610} \\
\bottomrule
\end{tabular}
\end{table}

\subsection{Le terme analytique, tenu pour invariant}
\label{sec:whitebox}

Le terme calculé l'est par la fonction \code{pinocchio.rnea} appliquée au modèle
construit directement depuis l'URDF. Trois décisions de conception en
découlent.

\parag{Un modèle réduit à sept axes.} L'URDF du Franka déclare neuf
articulations: les sept axes du bras et les deux glissières des doigts de la
pince. La boucle de commande ne pilote et ne mesure que les sept axes du bras.
Le modèle Pinocchio est donc réduit par \code{buildReducedModel}, qui verrouille
les deux glissières en position neutre et ramène la dimension de configuration
de 9 à 7. Sans cette réduction, tout appel RNEA en ligne échoue sur une erreur
de dimension; ce défaut était latent et n'est apparu qu'au premier appel réel
en boucle fermée.

\parag{Un terme analytique jamais modifié.} Le module RNEA est traité comme
une invariante du projet: aucun de ses paramètres n'est rendu optimisable, et
aucune de ses sorties n'est corrigée avant composition. C'est le point de
divergence délibéré avec DiffNEA~\cite{sutanto2020} et DeLaN~\cite{lutter2019iclr}.
Rendre l'URDF apprenable, comme le fait DiffNEA, détruirait la propriété sur
laquelle repose la contribution~1 — un modèle valide obtenu automatiquement
pour n'importe quelle description correcte — au bénéfice d'une précision
supplémentaire dont le résidu appris se charge par ailleurs. Le bénéfice de ce
gel est direct: la validité physique du terme dominant ne dépend en rien de la
qualité de l'apprentissage.

\parag{L'injection analytique de la charge.} Avant chaque évaluation de
RNEA, la masse portée $\delta$ est ajoutée à la masse du corps portant
l'effecteur, puis restaurée après l'appel. La partie \emph{connue} de l'effet de
la charge — surcroît de gravité, d'inertie et de couplage inertiel — est ainsi
traitée analytiquement, et seule la partie non modélisée reste à la charge du
résidu. Le mécanisme est identique à celui que Li \emph{et al.}~\cite{li2025ral}
présentent comme leur apport principal, implémenté ici indépendamment et
antérieurement; il constitue la première des deux voies de prise en compte de
la charge. En revanche, l'identification \emph{en ligne} de la charge par les
courants moteurs, qui fait le reste de leur contribution, a été rejetée: le
Franka est équipé de réducteurs, ce qui rend la relation courant-couple non
linéaire, contrairement aux actionneurs quasi directs de leur plateforme.

On notera dans la suite
%
\begin{equation}
    \taurbd(\q,\qd,\qdd,\delta) = \mathrm{RNEA}\bigl(\q,\qd,\qdd~;~
    \mathrm{URDF},\ \delta\bigr) \in \Rset^{7}
    \label{eq:taurbd}
\end{equation}
%
la sortie de ce terme analytique.

\subsection{Le résidu appris et ses trois degrés de garantie}
\label{sec:residu}

Reste à construire $\taures$. Plutôt que de présenter directement
l'architecture finale, il est plus informatif de la présenter comme le dernier
terme d'une progression en trois degrés de garantie, chacun apportant ce que le
précédent ne peut pas fournir. Cette manière de faire est celle de
Sutanto \emph{et al.}~\cite{sutanto2020}, qui comparent explicitement des
re-paramétrisations de contrainte croissante; elle rend visible ce que chaque
niveau coûte et ce qu'il achète.

\subsubsection{Degré~0: un résidu libre, régulier, conditionné par la charge}
\label{sec:degre0}

\parag{Encodage de l'état.} Le réseau résiduel ne reçoit pas directement
$\q$. Les positions articulaires sont d'abord transformées en leurs
représentations trigonométriques, ce qui élimine la discontinuité de
$\pm 2\pi$ et garantit la continuité spatiale de la fonction apprise. L'entrée
du réseau est le vecteur
%
\begin{equation}
    x = \bigl[\sin(\q),\ \cos(\q),\ \qd,\ \delta\bigr] \in \Rset^{22},
    \label{eq:encodage}
\end{equation}
%
de dimension $3\times 7 + 1 = 22$. La présence explicite de $\delta$ dans $x$
constitue la seconde voie de prise en compte de la charge: le résidu peut
apprendre des effets dépendant de la charge — frottements chargés, compliance de
transmission sous couple élevé — qu'un modèle rigide ne représente pas, même
avec l'inertie corrigée. C'est la combinaison des deux voies, analytique et
apprise, qui distingue ce travail des approches de la
section~\ref{sec:sota-boucle}, lesquelles n'en retiennent qu'une.

L'accélération $\qdd$ est délibérément absente de~\eqref{eq:encodage}. Elle
n'est nécessaire qu'à RNEA, où elle intervient hors ligne ou sous sa forme
désirée. Le résidu, lui, ne dépend que de la configuration, de la vitesse et de
la charge.

La présence de $\sin(\q), \cos(\q)$ dans cette entrée mérite une justification
propre, car un modèle de frottement sec et visqueux classique ne dépend que de
la vitesse (et, ici, de la charge) et non de la position. Ce n'est cependant
pas le modèle physique pertinent pour un réducteur harmonique comme ceux du
Franka~: l'engrènement des dents d'un réducteur harmonique produit une
ondulation de couple (\emph{torque ripple}) et une composante de frottement
répétable en fonction de l'angle articulaire, en plus de la composante
classique dépendant de la vitesse. C'est exactement ce qu'établissent
Hu \emph{et al.}~\cite{hu2026} pour l'erreur induite par le jeu mécanique
(\emph{backlash}) sur des robots industriels~: elle dépend de la direction de
la vitesse \emph{et} de la position articulaire, mais non de la charge portée.
La frontière tracée en~\eqref{eq:decompo} — charge à l'analytique, état de
mouvement (position et vitesse) au résidu — reprend directement cette
décomposition. Inclure $\q$ dans $x$ n'est donc pas une confusion entre
frottement et position~: c'est la manière dont ce travail représente une
composante de frottement räelle, répétable et dépendante de l'angle, propre
aux transmissions à réducteur harmonique, distincte de la composante
classique en $\qd$ que porte également le résidu.

\parag{Architecture et régularité.} Le résidu principal est un perceptron
multicouche de 4~couches cachées de 256 unités, sans activation sur la couche de
sortie, produisant $\taumlp \in \Rset^{7}$. Le choix de la fonction
d'activation est contraint: seules Mish et Softplus sont autorisées, et ReLU
est explicitement interdite dans tout le projet. La raison n'est pas la
performance d'apprentissage mais la physique de l'actionneur: la sortie du
réseau est un couple envoyé à des moteurs, et une activation non lisse produit
des discontinuités de dérivée qui se traduisent par des à-coups de couple. Le
même argument de régularité est invoqué par Liu \emph{et al.}~\cite{liu2024}
pour leur réseau d'énergie, où la nécessité de dériver deux fois impose des
activations lisses; il est ici transposé en argument de sécurité moteur. C'est
aussi le critère qui a conduit à rejeter deux mécanismes attractifs de la
littérature: les couches de Kolmogorov-Arnold recensées par
Toscano \emph{et al.}~\cite{toscano2025} et évaluées par
Wang \emph{et al.}~\cite{wang2025spel}, et le modèle de frottement de Stribeck
paramétrique de Lutter \emph{et al.}~\cite{lutter2019iros}, qui repose sur un
terme en $\mathrm{sign}(\qd)$ non lisse.

Le critère de régularité $C^{1}$ ne distingue cependant pas, à lui seul, Mish
de deux autres activations lisses courantes, GELU et Tanh, et ce choix précis
mérite d'être justifié séparément — étant entendu que Liu \emph{et al.}~
\cite{liu2024}, cités ci-dessus pour l'argument de régularité, n'utilisent pas
Mish dans leur propre réseau (ils emploient Softplus en sortie pour la
non-négativité) et ne fournissent donc aucun précédent direct à invoquer ici.
Tanh est écartée pour une raison de saturation~: bornée dans $(-1,1)$, elle
produit un gradient qui s'annule pour les grandes activations, ce qui
ralentirait l'apprentissage des régimes de couple élevé — précisément ceux que
ce travail cherche à bien représenter sur les axes porteurs. Entre GELU et
Mish, toutes deux non bornées supérieurement, bornées inférieurement et non
monotones au voisinage de zéro (propriétés qui leur valent d'être préférées à
Tanh dans la littérature récente sur les réseaux profonds), le choix de Mish
retenu ici repose sur ces propriétés génériques de la fonction elle-même — en
particulier l'auto-régularisation de sa composante non monotone au voisinage
de zéro, rapportée comme bénéfique à l'entraînement dans sa littérature
d'origine — et non sur une comparaison empirique menée dans ce projet. Aucune
étude de sensibilité Mish contre GELU n'a été réalisée ici~: ce n'est pas un
axe de la matrice d'ablation implémentée (section~\ref{sec:limites}), et le
choix doit donc être lu comme un choix de conception par défaut, documenté et
assumé, plutôt que comme un résultat comparatif.

Ce degré~0 ne garantit rien. Un tel résidu peut prédire des couples hors des
limites d'actionneur et, plus grave, injecter de l'énergie dans le système —
c'est-à-dire se comporter comme une source alors qu'il est censé représenter des
pertes.

\subsubsection{Degré~1: des garanties imposées par la fonction objectif}
\label{sec:degre1}

Le premier remède consiste à contraindre l'optimisation. L'entraînement minimise
une perte composite formée d'un terme de fidélité aux données et d'un lagrangien
augmenté portant deux contraintes d'inégalité. Pour un lot de $N$ échantillons
indexés par $i$,
%
\begin{equation}
    \mathcal{L}_{\text{données}}
    = \frac{1}{N}\sum_{i=1}^{N}
      \bigl\| \taupred^{(i)} - \taureal^{(i)} \bigr\|_2^{2} .
    \label{eq:Ldata}
\end{equation}
%
Les deux contraintes sont
%
\begin{align}
    g_{\mathrm{couple}}^{(i)} &= \bigl|\taupred^{(i)}\bigr| - \bar{\tau} \le 0
    \qquad \text{(par axe)}, \label{eq:gtorque}\\
    g_{\mathrm{dis}}^{(i)} &= \bigl(\taures^{(i)}\bigr)^{\top}\qd^{(i)} \le 0 ,
    \label{eq:gdissip}
\end{align}
%
où $\bar{\tau}$ est le vecteur des limites de couple du
tableau~\ref{tab:constantes}. Elles ne jouent pas le même rôle:
\eqref{eq:gtorque} est une contrainte de \emph{sécurité}, qui interdit au modèle
de prédire un couple irréalisable, alors que~\eqref{eq:gdissip} est la
contrainte \emph{informative}, qui interdit au résidu total d'injecter de
l'énergie.

En posant $\psi_\bullet = \max(0, g_\bullet)$ la violation positive, la perte
totale s'écrit
%
\begin{equation}
    \mathcal{L}
    = \mathcal{L}_{\text{données}}
    + \underbrace{\overline{\lambda_{t}\,\psi_{\mathrm{couple}}
      + \tfrac{\rho}{2}\,\psi_{\mathrm{couple}}^{2}}
    + \overline{\lambda_{d}\,\psi_{\mathrm{dis}}
      + \tfrac{\rho}{2}\,\psi_{\mathrm{dis}}^{2}}}
      _{\text{lagrangien augmenté}} ,
    \label{eq:Ltotal}
\end{equation}
%
où la barre désigne la moyenne sur le lot, $\rho = 1$ est le poids de
pénalisation quadratique, et $\lambda_{t}, \lambda_{d} \ge 0$ sont les
multiplicateurs de Lagrange, mis à jour une fois par époque par montée duale.
La valeur $\rho = 1$ elle-même est fixée arbitrairement, à dire d'expérience,
et n'a fait l'objet d'aucune étude de sensibilité ni d'aucune règle de mise à
jour adaptative dans ce travail — une lacune reconnue ici plutôt que
découverte a posteriori, alors même que l'instabilité de la montée duale est
classiquement associée à un mauvais calibrage de ce paramètre. La
section~\ref{sec:res-physique} montrera que la montée duale converge
proprement sur données saines avec cette valeur, ce qui atténue mais
n'élimine pas ce risque~: rien ne garantit que $\rho = 1$ reste un bon choix
hors du régime testé ici, et une étude de sensibilité sur $\rho$ fait partie
des axes d'ablation non couverts (section~\ref{sec:limites}).
%
\begin{equation}
    \lambda_\bullet \leftarrow \max\bigl(0,\ \lambda_\bullet
    + \rho\,\overline{\psi_\bullet}\bigr) .
    \label{eq:dual}
\end{equation}
%
Ce mécanisme, repris de Djeumou \emph{et al.}~\cite{djeumou2022}, présente
l'avantage décisif sur une pénalisation à poids fixe de ne pas exiger le réglage
manuel d'un compromis: par~\eqref{eq:dual}, le multiplicateur croît tant que la
contrainte est violée et se relâche dès qu'elle est satisfaite. C'est aussi
pourquoi le
programme d'annelage des contraintes proposé par
Prabhakar \emph{et al.}~\cite{prabhakar2026} a été écarté: la montée duale
réalise déjà, de manière adaptative, ce que cet annelage impose de manière
prescrite.

Ce degré~1 apporte beaucoup, mais sa garantie reste d'une nature faible. Elle
est \emph{statistique et rétrospective}: les contraintes sont évaluées sur les
échantillons vus à l'entraînement, et rien n'assure qu'un état non visité les
respecte. La section~\ref{sec:res-physique} montrera d'ailleurs que la montée
duale, si elle converge proprement sur des données saines, peut se déstabiliser
lorsque les données ne le sont pas.

\subsubsection{Degré~2: une garantie obtenue par construction}
\label{sec:frictionnet}

Le troisième degré consiste à retirer la partie dissipative du résidu du domaine
de ce qui est \emph{espéré} pour la placer dans celui de ce qui est \emph{vrai}.
Un sous-module dédié, appelé ici module de frottement, porte cette partie. Son
architecture est un perceptron $22 \rightarrow 64 \rightarrow 64 \rightarrow 7$
à activations Mish, dont la sortie $d \in \Rset^{7}$ est transformée en une
matrice de frottement diagonale
%
\begin{equation}
    D(x) = \mathrm{diag}\bigl(\mathrm{Softplus}(d) + \varepsilon\bigr),
    \qquad \varepsilon = 10^{-6},
    \label{eq:Dmat}
\end{equation}
%
strictement définie positive, puis en un couple de frottement
%
\begin{equation}
    \taufric = -D(x)\,\qd .
    \label{eq:taufric}
\end{equation}
%
La puissance dissipée vérifie alors, pour tout état et tout jeu de poids,
%
\begin{equation}
    \taufric^{\top}\qd
    = -\sum_{j=1}^{7}\bigl(\mathrm{Softplus}(d_j)+\varepsilon\bigr)\,\qd_j^{2}
    \le 0 ,
    \label{eq:dissip}
\end{equation}
%
c'est-à-dire que le couple~\eqref{eq:taufric}, construit depuis la matrice
strictement définie positive~\eqref{eq:Dmat}, retire toujours de l'énergie au
système. Contrairement au degré~1, la propriété~\eqref{eq:dissip} ne dépend ni
des données ni de l'entraînement: c'est une identité algébrique.

\parag{Pourquoi une matrice diagonale.} Le motif de factorisation
garantissant la définie-positivité est celui de
Lutter \emph{et al.}~\cite{lutter2019iclr}, appliqué à la dissipation par
Liu \emph{et al.}~\cite{liu2024} et Duong \emph{et al.}~\cite{duong2024}. Une
factorisation de Cholesky complète d'une matrice $7\times 7$ demanderait
28~entrées triangulaires inférieures. L'argument géométrique de
Wang \emph{et al.}~\cite{wang2025spel} établit que des liaisons rotoïdes
tournant chacune autour d'un axe unique produisent une dissipation de rang~1 par
corps, qui se réduit en espace articulaire à un scalaire par axe. Cette
hypothèse a été vérifiée spécifiquement sur le modèle du Franka par une analyse
dédiée du modèle Pinocchio: les sept articulations sont bien du type révolute à
axe unique, et les coefficients de frottement et d'amortissement déclarés dans
l'URDF sont tous nuls, donc tous à apprendre. La forme diagonale est ainsi
physiquement justifiée et réduit de 28 à 7 le nombre de paramètres de
dissipation, soit 80\% de moins, pour un total de 6\,087 paramètres dans le
sous-module.

Le résidu total composé est finalement
%
\begin{equation}
    \taures(\q,\qd,\delta) = \taumlp(x) + \taufric(x),
    \label{eq:greybox}
\end{equation}
%
et, injecté dans~\eqref{eq:decompo}, il donne l'architecture représentée sur la
figure~\ref{fig:archi}. On notera que la contrainte~\eqref{eq:gdissip} du
degré~1 n'est pas rendue inutile par le degré~2: elle étend à $\taumlp$, dont
la sortie est libre, la propriété que~\eqref{eq:dissip} garantit pour
$\taufric$ seul. Les deux degrés sont complémentaires, l'un couvrant ce que
l'autre laisse libre.

\begin{figure}[htbp]
\centering
\begin{tikzpicture}[
    font=\small,
    box/.style={draw, rounded corners=2pt, minimum height=8mm, align=center,
                inner sep=3pt},
    white/.style={box, fill=black!4},
    learn/.style={box, fill=black!12},
    ar/.style={-{Latex[length=2mm]}, thick},
    sum/.style={draw, circle, inner sep=1pt}
]
% voie analytique (rangée haute)
\node[white, text width=36mm] (rnea) at (5.2,3.0)
    {\textbf{terme calculé}\\RNEA (URDF, $\delta$)\\
     $\taurbd$ \eqref{eq:taurbd}};
% voie apprise (rangée basse)
\node[box, text width=22mm] (state) at (1.3,0)
    {$\q,\qd,\qdd,\delta$};
\node[box, text width=30mm] (enc) at (5.0,0)
    {encodage \eqref{eq:encodage}\\$x \in \Rset^{22}$};
\node[learn, text width=32mm] (mlp) at (9.6,1.3)
    {MLP $4\times256$, Mish\\$\taumlp$ (degré~0)};
\node[learn, text width=30mm] (fric) at (9.6,-1.5)
    {module de frottement\\MLP 22-64-64-7\\$\taufric=-D(x)\,\qd$ (degré~2)};
\node[sum] (s1) at (12.9,0) {$+$};
\node[sum] (s2) at (12.9,3.0) {$+$};
\node[box] (out) at (14.3,3.0) {$\taupred$};

\draw[ar] (state.north) |- (rnea.west);
\draw[ar] (state.east) -- (enc.west);
\draw[ar] (enc.east) -- (mlp.west);
\draw[ar] (enc.east) -- (fric.west);
\draw[ar] (mlp.east) -- (12.3,1.3) |- (s1.west);
\draw[ar] (fric.east) -- (12.3,-1.5) |- (s1.west);
\draw[ar] (s1.north) -- node[right, pos=0.5] {$\taures$} (s2.south);
\draw[ar] (rnea.east) -- (s2.west);
\draw[ar] (s2.east) -- (out.west);
\end{tikzpicture}
\caption{Architecture grise résultant de~\eqref{eq:decompo}
et~\eqref{eq:greybox}. Le terme calculé RNEA, jamais appris, porte la dynamique
du corps rigide et la part connue de la charge; les deux blocs appris (gris
foncé) portent le résidu. Le module de frottement est structurellement
dissipatif par~\eqref{eq:dissip} (degré~2), tandis que la sortie libre du MLP
(degré~0) n'est contrainte que par la perte~\eqref{eq:Ltotal} (degré~1). Tous
les blocs appris reçoivent la charge $\delta$ en entrée.}
\label{fig:archi}
\end{figure}

\subsection{Rendre le résidu identifiable: excitation et hygiène des données}
\label{sec:donnees}

Un résidu contraint n'est utile que s'il est identifiable, ce qui exige des
données couvrant l'espace d'état atteignable. La stratégie retenue est en deux
niveaux, puis suivie d'un filtrage dont la justification constitue l'un des
résultats du travail.

\parag{Choix des configurations centrales.} Dix configurations articulaires
$\q_{\mathrm{c}}$ sont tirées par une séquence de Sobol brouillée en
dimension~7, bornée par les limites articulaires de l'URDF. Un échantillonnage
quasi aléatoire à faible discrépance couvre l'espace plus uniformément qu'un
tirage uniforme indépendant, ce qui importe lorsque le nombre de
configurations est faible devant la dimension — mais dix points, même bien
répartis par une séquence de Sobol, restent un échantillonnage extrêmement
clairsemé d'un espace à 7~dimensions~: la distance entre plus proches voisins
y est grande, et de larges régions de l'espace articulaire ne sont couvertes
par aucune configuration centrale, seulement par l'excursion locale, de faible
amplitude (0,05 à 0,30~rad, voir~\eqref{eq:fourier}), des trajectoires
d'excitation qui les entourent. Ce choix est un compromis budgétaire assumé
plutôt qu'une couverture revendiquée de l'espace atteignable~: chaque
configuration centrale coûte 5\,000 pas de simulation par charge, et
l'élargir à un nombre de points croissant avec la dimension (ce que
suggérerait une couverture uniforme au sens strict) était hors du budget de
calcul de ce travail. La conséquence directe est que le modèle appris n'est
identifié, et donc valide, que dans un voisinage de ces dix configurations et
de leur excursion locale — c'est une limite reconnue explicitement en
section~\ref{sec:limites}, et non un défaut découvert après coup.

\parag{Trajectoires d'excitation.} Autour de chaque configuration centrale,
une série de Fourier à 5~harmoniques est générée indépendamment pour chaque
axe:
%
\begin{equation}
    \q_j(t) = \q_{\mathrm{c},j}
    + \sum_{k=1}^{5} A_{kj}\,\sin\!\bigl(\omega_{kj}\,t + \varphi_{kj}\bigr),
    \label{eq:fourier}
\end{equation}
%
avec $A_{kj}$ tiré uniformément dans $[0{,}05~;~0{,}30]$~rad,
$\omega_{kj}$ dans $[0{,}50~;~3{,}00]$~rad$\cdot$s$^{-1}$ et $\varphi_{kj}$
dans $[0~;~2\pi]$. La série de Fourier finie est le choix standard pour
l'excitation en identification de dynamique~\cite{deng2024}; son avantage
propre ici est que les vitesses et accélérations s'obtiennent par dérivation
analytique de~\eqref{eq:fourier}, ce qui fournit un $\qdd$ exact — et non
différencié numériquement — pour l'évaluation de RNEA. Chaque segment dure 5~s
échantillonnées à \SI{1}{kHz}, soit 5\,000 pas, et dix segments sont exécutés par
valeur de charge, ce qui porte le jeu nominal à 150\,000 échantillons pour les
trois charges.

\parag{Source du couple mesuré.} Les trajectoires sont rejouées dans le
moteur physique d'Isaac~Sim, dont le modèle du Franka porte des frottements, un
amortissement et une compliance d'actionneur propres; le couple mesuré
$\taureal$ est lu directement sur les articulations. Le résidu cible
$\taureal - \taurbd$ reflète donc de la dynamique non modélisée d'un moteur
physique indépendant, et non un modèle de frottement synthétique. Un générateur
antérieur, fondé sur un résidu de Coulomb et visqueux réglé à la main, a été
conservé mais uniquement comme test de bon fonctionnement de la chaîne sur
processeur: il n'a pas de valeur scientifique, puisque le réseau y réapprend un
modèle inventé.

\parag{Filtrage de saturation.} Les échantillons dont le couple mesuré
atteint la limite de l'actionneur sont écartés par le critère
%
\begin{equation}
    \bigl|\tau_{\mathrm{real},j}\bigr| \le 0{,}97\,\bar{\tau}_j
    \quad \text{pour tout } j ,
    \label{eq:satfilter}
\end{equation}
%
c'est-à-dire une marge de 3\% sous la limite. Le raisonnement, développé et
quantifié en section~\ref{sec:res-satfix}, est le suivant: un échantillon saturé
porte comme couple mesuré la valeur d'écrêtage de l'actionneur et non la
dynamique réelle, alors que $\taurbd$ est calculé depuis l'accélération \emph{de
référence}, supposée atteinte; le résidu cible correspondant est alors dénué de
sens physique. Un filtre naïf $|\taureal| > \bar{\tau}$ ne les capture pas,
puisqu'un échantillon écrêté se situe exactement \emph{à} la limite. Après
filtrage, le jeu de données retient 148\,304~échantillons sur les 150\,000
nominaux — soit un peu plus de 1\% écartés — répartis sur trois valeurs de
charge (\SIlist{0;1;3}{kg}) et séparés en 90\% d'entraînement et 10\% de validation.

\subsection{Fermer la boucle sur le modèle appris}
\label{sec:boucle}

\subsubsection{Loi de commande}

La loi de commande est un couple calculé augmenté d'un correcteur
proportionnel-dérivé. À chaque pas, pour un état mesuré $(\q,\qd)$ et une
consigne $(\q_{\mathrm{d}},\qd_{\mathrm{d}},\qdd_{\mathrm{d}})$ issue de la
trajectoire planifiée,
%
\begin{equation}
    \taucmd = \underbrace{\taurbd(\q,\qd,\qdd_{\mathrm{d}},\delta)
    + \taures(\q,\qd,\delta)}_{\text{anticipation, modèle appris}}
    + \underbrace{K_p\,e + K_d\,\dot{e}}_{\text{rétroaction}},
    \qquad
    \begin{aligned}
        e &= \q_{\mathrm{d}} - \q, \\
        \dot{e} &= \qd_{\mathrm{d}} - \qd,
    \end{aligned}
    \label{eq:ctpd}
\end{equation}
%
la sortie étant finalement écrêtée axe par axe aux limites du
tableau~\ref{tab:constantes}:
$\taucmd \leftarrow \mathrm{clip}(\taucmd, -\bar{\tau}, +\bar{\tau})$. Cet
écrêtage est une couche de sécurité indépendante de l'apprentissage: même si le
modèle prédisait un couple irréalisable, aucune consigne hors limites ne peut
être publiée. On soulignera que $\qdd_{\mathrm{d}}$ provient du planificateur et
n'est jamais une accélération différenciée depuis les capteurs, ce qui referme
l'argument ouvert en section~\ref{sec:sota-archi}.

\subsubsection{Des gains dérivés d'une borne mesurée}
\label{sec:gains}

Les gains ne sont pas réglés à la main. Le résultat de
Liu \emph{et al.}~\cite{liu2024} établit que, pour une loi de la
forme~\eqref{eq:ctpd} avec une erreur de modèle \emph{strictement bornée au
sens du pire cas} par $\epsilon_j$ sur l'axe $j$, l'erreur de suivi converge
dès lors que les gains dominent cette borne avec une marge suffisante. Il faut
préciser d'emblée un point qui conditionne la portée de tout ce qui suit~: la
démonstration exige une borne valable pour \emph{tout} état, et non une borne
statistique. La section~\ref{sec:res-bound} montrera que la borne effectivement
injectée dans~\eqref{eq:gains} ci-dessous est un quantile mesuré à 99,9~\% et
non le maximum absolu observé~; ce choix, justifié et quantifié à cet endroit,
affaiblit la garantie théorique de ce paragraphe~: la certification au sens de
Lyapunov qui suit n'est donc, en toute rigueur, valable que pour la fraction
des conditions de fonctionnement couvertes par ce quantile, le reste étant
couvert non par la preuve mais par l'écrêtage matériel indépendant décrit plus
haut. On en déduit, en imposant de surcroît une condition d'amortissement
critique qui évite tout dépassement,
%
\begin{equation}
    K_{d,jj} = \kappa\,\epsilon_j , \qquad
    K_{p,jj} = \frac{K_{d,jj}^{2}}{4} ,
    \label{eq:gains}
\end{equation}
%
où $\kappa \ge 1$ est une marge de sécurité multiplicative, fixée à
$\kappa = 2$, et $K_p, K_d$ sont diagonales définies positives. Le point
méthodologique est que $\epsilon_j$ n'est pas postulé: il est mesuré sur la
partition de validation du modèle entraîné, par un outil dédié qui reconstruit
exactement la même partition que l'entraînement et rapporte, par axe, l'erreur
quadratique moyenne et les quantiles de $|\taupred - \taureal|$. Les gains sont
donc une conséquence chiffrée de la qualité du modèle, et non un réglage
indépendant. Les valeurs obtenues et la discussion du choix du quantile figurent
en section~\ref{sec:res-bound}.

\subsubsection{Environnement d'exécution et contraintes de réalisation}
\label{sec:exec}

\parag{Le planificateur et le matériel simulé.} La planification est assurée
par MoveIt2 sous ROS2~Jazzy, dans sa configuration standard, avec résolution de
cinématique inverse par le solveur KDL. Ce choix est délibérément
minimaliste: la planification n'est pas l'objet de la contribution, et lui
substituer un planificateur de recherche — par exemple un planificateur
neuronal résolvant l'équation d'Eikonal comme celui de
Ni et Qureshi~\cite{ni2024} — aurait introduit une variable de confusion
supplémentaire dans un système déjà difficile à déboguer. Le matériel simulé, en
revanche, a changé en cours de projet, et la justification de ce changement fait
partie de la méthode. La première version utilisait Isaac~Sim comme robot
simulé, via l'espace de travail ROS2 fourni par l'éditeur, ce qui imposait une
dépendance de compilation externe au dépôt et non versionnée. Elle a été
remplacée par MuJoCo à travers le paquet \code{mujoco_ros2_control}, installable
par le gestionnaire de paquets de la distribution. Deux critères ont motivé la
bascule: la suppression de la dépendance externe, et surtout un avantage
structurel — le modèle MJCF utilisé déclare un actionneur de couple direct
unique par axe, sans asservissement de position interne concurrent, alors que le
modèle USD d'Isaac~Sim comportait un asservissement de position raide qu'il
fallait neutraliser explicitement pour que les consignes de couple aient une
autorité. Cette classe de défaut disparaît par construction. Deux simulateurs
indépendants coexistent donc dans ce travail~: Isaac~Sim reste la source hors
ligne des données d'entraînement, MuJoCo est le \og matériel simulé \fg{}
piloté en temps réel. Cet écart de domaine entre simulateur d'entraînement et
simulateur d'exécution — un \emph{sim-to-sim gap} — n'est pas anodin et doit
être explicité plutôt que seulement mentionné~: les deux moteurs physiques
n'implémentent ni le même solveur de contact, ni le même modèle de frottement
d'actionneur, ni nécessairement les mêmes paramètres numériques d'intégration,
de sorte que le couple $\taureal$ appris depuis Isaac~Sim n'est pas
garanti représenter fidèlement la dynamique que le modèle rencontre effectivement
une fois inséré dans la boucle de commande sous MuJoCo. Ce travail ne
neutralise pas cet écart~: il l'assume comme une limite explicite (voir
section~\ref{sec:limites}), et le distingue soigneusement du sim-to-real
gap visé par la contribution~4~: la validation en boucle fermée de
la section~\ref{sec:res-loop} démontre que le modèle appris pilote
\emph{effectivement} un simulateur physique \emph{différent} de celui sur
lequel il a été identifié — un test de robustesse involontaire mais réel — et
non qu'il piloterait un robot réel. Le prix précis de ce choix, et pourquoi il
a néanmoins été retenu (suppression d'une dépendance de compilation externe et
d'un asservissement de position concurrent, voir ci-dessus), est discuté plus
loin en section~\ref{sec:limites}.

\parag{Réalisation temps réel.} La boucle est réalisée par un n\oe ud ROS2
dédié, cadencé par une minuterie à \SI{1000}{Hz}. À chaque tic: lecture du dernier
état articulaire reçu, interpolation linéaire de la consigne dans la trajectoire
courante, appel de la loi~\eqref{eq:ctpd}, publication du vecteur de couples sur
le contrôleur d'effort de \code{ros2_control}. Conformément à l'avertissement
d'Ibarz \emph{et al.}~\cite{ibarz2021} sur l'exécution asynchrone, trois détails
de réalisation se sont révélés critiques et méritent d'être documentés comme
éléments de méthode plutôt que comme anecdotes.

\begin{description}
    \item[Indexation par nom, jamais par position.] Le diffuseur d'état
    articulaire de \code{ros2_control} publie \emph{toutes} les articulations
    déclarées auprès du gestionnaire de contrôleurs, triées par ordre
    alphabétique, y compris les deux glissières de la pince. Extraire les sept
    premiers éléments du message donne donc un vecteur mélangé: les deux
    doigts précèdent \code{panda_joint1} et rejettent \code{panda_joint6} et
    \code{panda_joint7} hors des sept premières positions. Le vecteur de
    configuration passé à RNEA et au réseau était par conséquent incohérent à
    chaque pas, avec pour symptôme la chute systématique du bras à chaque
    bascule vers le mode couple. La règle retenue — indexer par nom — a dû
    être appliquée deux fois, à deux endroits indépendants du code
    (souscription à l'état articulaire, puis interpolateur de trajectoire).
    \item[Fenêtre de péremption des consignes.] La trajectoire reçue est
    considérée comme périmée au bout de 2~s de temps mural après réception,
    après quoi le n\oe ud se replie sur une compensation de gravité. Ce
    dispositif de sûreté a une conséquence non évidente sur les couches
    supérieures: tout producteur de trajectoire doit soit rester sous cette
    durée, soit republier périodiquement pour maintenir la consigne fraîche.
    \item[Cadence de la boucle intergicielle.] La configuration de
    démonstration héritée cadençait la boucle de l'intergiciel à \SI{100}{Hz}, ce
    qui produisait une oscillation visible en maintien de position. Le passage
    à \SI{1000}{Hz}, conforme à l'objectif, l'a supprimée.
\end{description}

\subsection{Transférer sous contrainte: affinage à ossature gelée}
\label{sec:finetune}

Le protocole de transfert vers un robot physique est en deux temps. Le modèle
est d'abord pré-entraîné sur la totalité des données de simulation. Puis, pour
l'adaptation, toutes les couches sont gelées sauf les deux dernières couches
linéaires, et l'optimisation reprend sur un petit enregistrement réel de
mouvements exploratoires, pendant un nombre volontairement faible de pas de
gradient (100 par défaut, plage 50 à 200) avec un pas d'apprentissage réduit à
$10^{-4}$:
%
\begin{equation}
    \theta^{*} = \arg\min_{\theta \in \Theta_{\mathrm{libre}}}
    \mathcal{L}\bigl(\theta,\ \mathcal{D}_{\text{réel}}\bigr) .
    \label{eq:finetune}
\end{equation}
%
Deux critères motivent ce dispositif. D'abord, le gel de l'ossature limite
l'oubli catastrophique: la représentation apprise sur $1{,}5\cdot 10^{5}$
échantillons ne doit pas être détruite par quelques milliers d'échantillons
réels. Ensuite, et c'est le point spécifique à une approche informée par la
physique, l'affinage se déroule sous la perte complète~\eqref{eq:Ltotal},
contraintes incluses: le réseau ne peut pas s'ajuster au bruit de mesure en
violant la dissipativité ou les limites de couple. Les contraintes jouent donc
le rôle d'un régularisateur physique pendant la phase où le risque de
surapprentissage est maximal.

Cet argument appelle une précision importante sur son domaine de validité~:
il concerne le bruit de mesure d'un \emph{capteur physique réel} — jeu de
codeur, quantification, vibration — sur $\mathcal{D}_{\text{réel}}$, un
enregistrement effectué sur le Franka matériel. Il ne s'applique pas tel quel
si $\mathcal{D}_{\text{réel}}$ venait à être approché, pour test du
mécanisme, par un enregistrement supplémentaire issu d'Isaac~Sim~: un
simulateur ne génère pas nativement ce bruit de capteur, et l'argument de
protection contre le surapprentissage du bruit y perdrait sa justification
physique, même si les contraintes resteraient utiles comme régularisateur
générique contre le surapprentissage d'un petit échantillon. Comme le rappelle
la section~\ref{sec:limites}, ce protocole n'a en pratique jamais été exécuté
sur un enregistrement réel dans ce travail~: l'argument ci-dessus est donc, à
ce stade, une justification de conception pour le cas d'usage visé, non un
bénéfice mesuré. Le protocole de gel est repris de
Duong \emph{et al.}~\cite{duong2024}; son exécution sous perte contrainte en
est l'extension propre à ce travail.

\subsection{Une tâche de préhension comme test d'intégration}
\label{sec:tache}

La dernière pièce n'apporte aucune nouveauté scientifique: son rôle est
d'exercer simultanément toutes les précédentes sur une tâche complète, et donc
de constituer un test d'intégration au sens fort. La scène MuJoCo comporte un
sol, un éclairage, et un cube libre de 4~cm de côté posé à 50~cm devant la base
du robot, largement dans l'espace de travail (portée d'environ 85{,}5~cm). La
séquence est réalisée par une machine à états dont les phases sont: repos,
pré-approche, approche, fermeture, levée, maintien, relâchement, et un état
d'abandon. Le mouvement du bras est délégué à un client qui construit une cible
cartésienne (contrainte de position et contrainte d'orientation), demande un
plan à MoveIt2, transmet la trajectoire obtenue au n\oe ud de commande, puis
vérifie la convergence. Trois choix de conception, tous issus de
l'expérimentation en boucle réelle, méritent d'être justifiés.

\parag{Un critère de convergence cartésien plutôt qu'articulaire.} Le premier
critère utilisé était un seuil uniforme en radians sur l'erreur articulaire
maximale. Il est physiquement mal posé: une même erreur en radians sur l'axe~1
ou sur l'axe~7 ne produit pas du tout le même écart de l'effecteur. Il a été
remplacé par une vérification de la distance cartésienne réelle de la bride,
calculée par cinématique directe avec Pinocchio sur le même URDF que le reste de
la chaîne. Deux tolérances distinctes sont utilisées: 2~cm pour la descente
finale, où la précision compte, et 4~cm pour les points de passage
intermédiaires de survol, où elle ne compte pas.

\parag{L'enregistrement des obstacles auprès du planificateur.} MoveIt2
ignorait initialement l'existence du cube: les plans étaient déclarés libres,
mais l'erreur de suivi réelle amenait le bras à toucher effectivement l'objet,
ce que le moteur physique traduisait par un contact absorbant le couple
commandé — symptôme visuellement identique à un axe bloqué. Le cube est
désormais enregistré comme obstacle de la scène de planification, avec une
taille majorée à 12~cm au lieu de ses 4~cm réels afin d'absorber l'erreur de
suivi, puis attaché à la pince juste avant la descente et détaché au
relâchement. Le sol a été enregistré de la même façon.

\parag{Le découpage du mouvement en points de passage.} Le trajet depuis la
position de repos vers la pose de pré-approche est interpolé en trois
sous-mouvements par interpolation complète en $SE(3)$ (position linéaire et
orientation par interpolation sphérique), plutôt qu'exécuté en un seul grand
mouvement. Le critère est ici la robustesse du suivi: un grand débattement
demandé en une fois a produit des échecs de convergence que des sous-mouvements
plus courts n'ont pas.

% =====================================================================
\section{Validation et résultats}
\label{sec:resultats}
% =====================================================================

\subsection{Ce qui doit être démontré, et par quelle grandeur}
\label{sec:res-quoi}

Aucune mesure unique ne peut valider l'ensemble de la méthode. La stratégie
adoptée consiste à associer à chaque contribution revendiquée une grandeur
susceptible de la démentir, puis à progresser du plus contrôlé au plus intégré.
Le choix de ces grandeurs demande d'être justifié avant d'être appliqué, car
deux d'entre elles sont moins évidentes qu'il n'y paraît.

\parag{L'erreur de couple doit être lue relativement à la limite d'axe.}
Une erreur de prédiction exprimée en newton-mètres n'est pas comparable d'un axe
à l'autre lorsque les limites d'actionneur diffèrent d'un facteur~7
(tableau~\ref{tab:constantes}). La grandeur qui compte pour la commande est la
fraction de l'autorité de l'actionneur consommée par l'erreur de modèle, soit
$\mathrm{RMSE}_j / \bar{\tau}_j$. Les deux lectures — absolue et relative — sont
rapportées, et la section~\ref{sec:res-stage1} montrera qu'elles ordonnent les
axes en sens inverse.

\parag{Une moyenne ne suffit pas à valider un modèle destiné à commander.}
L'avertissement de Deng \emph{et al.}~\cite{deng2024} est explicite: un modèle
peut afficher une erreur moyenne excellente tout en produisant des prédictions
localement absurdes que la moyenne dissimule. Les queues de distribution sont
donc rapportées séparément (section~\ref{sec:res-bound}), et ce sont elles, non
la moyenne, qui alimentent le calcul des gains~\eqref{eq:gains}. C'est
d'ailleurs par cette lecture que le défaut de données le plus important du
travail a été découvert.

Les cinq niveaux de vérification sont alors les suivants. \emph{Niveau~0} —
cohérence structurelle: chaque module possède un test vérifiant les dimensions,
l'absence d'activation interdite et, pour le module de frottement, la
propriété~\eqref{eq:dissip} sur des états tirés aléatoirement; c'est un cas
simple où la solution est connue analytiquement, puisque la puissance dissipée
\emph{doit} être négative quels que soient les poids. \emph{Niveau~1} —
précision hors ligne sur une partition de validation de 10\% jamais vue à
l'entraînement, décomposée par axe (section~\ref{sec:res-stage1}).
\emph{Niveau~2} — tenue des garanties physiques au fil des époques
(section~\ref{sec:res-physique}), puis une expérience contrôlée avec hypothèse
préalable, seul endroit du travail où une hypothèse a été confirmée par une
comparaison à réglages identiques (section~\ref{sec:res-satfix}).
\emph{Niveau~3} — boucle fermée: capture numérique de l'état articulaire mesuré
et de la décomposition du couple commandé pendant une exécution réelle, pour
vérifier que le modèle appris pilote effectivement le robot simulé
(section~\ref{sec:res-loop}). \emph{Niveau~4} — tâche complète: taux de
réussite et erreur finale de l'effecteur sur douze exécutions instrumentées
(section~\ref{sec:res-task}). Ce qui manque à ce dispositif est identifié
explicitement en section~\ref{sec:limites}: il n'y a pas de validation sur
robot physique, et la mesure de latence d'inférence qui substantierait la
contribution~2 n'a pas été réalisée.

\subsection{Précision du modèle appris}
\label{sec:res-stage1}

% TODO(hugo): compléter ici le modèle exact de GPU/CPU utilisé pour
% l'entraînement (machine `hci-student`) ainsi que le temps d'entraînement
% mural mesuré pour les 200 époques ci-dessous ; ni le dépôt du projet ni
% tracking/PROJECT_STATE.md ne précisent la référence matérielle exacte ou
% une durée chiffrée au moment de cette révision, seulement une mention
% générique de « machine GPU ». Ne pas inventer de valeur : demander la
% mesure exacte avant soutenance.
% TODO(siunitx): conversion \SI{}{} non faite pour les grandeurs numériques
% des grands tableaux de cette section (tab:perjoint, tab:exp, tab:tau,
% tab:stage4) et des figures associées, par souci de temps ; les unités
% (Nm, rad, m, %) y sont donc restées en notation manuelle « 0,80~Nm » etc.
% Cohérent en soi, mais un passage manuel de conversion vers \SI{}{}/\num{}
% reste à faire pour une cohérence totale avec le reste du document.
La configuration de référence, notée ici \code{isaac-satfix}, est un perceptron
de 4~couches de 256 unités à activation Mish, module de frottement activé,
optimiseur Adam à pas $10^{-3}$, lots de 256 échantillons, 200~époques,
$\rho = 1$, entraîné sur les 148\,304 échantillons issus d'Isaac~Sim aux trois
charges. Sa meilleure perte de validation vaut
$\mathcal{L}_{\mathrm{val}} = 0{,}3995$, atteinte à l'époque~144 et non à la
dernière époque. Elle se compare directement à celle d'une exécution
antérieure identique en tous points sauf le filtrage
de saturation~\eqref{eq:satfilter}, notée \code{isaac-first-real}, qui plafonne
à $1{,}4767$; cette comparaison est l'objet de la
section~\ref{sec:res-satfix}. Quatre exécutions plus anciennes, menées sur
résidu synthétique, ont servi uniquement à valider la mécanique de la chaîne sur
processeur et ne sont pas comparables aux précédentes, leur cible étant un
modèle de frottement inventé, plus facile à apprendre qu'une dynamique de
moteur physique; elles ne sont mentionnées que pour la traçabilité.

La figure~\ref{fig:courbes} montre les courbes de perte de validation des deux
exécutions sur données physiques. Deux observations quantitatives s'en dégagent.
D'abord, l'écart de niveau est net et constant: la courbe de l'exécution
corrigée s'établit autour de 0{,}42 tandis que celle de l'exécution initiale
plafonne autour de 1{,}6, soit un facteur voisin de~4. Ensuite, et c'est plus
instructif, la \emph{régularité} des deux courbes diffère: l'exécution initiale
présente des sursauts marqués (2{,}40 à l'époque~50 puis 2{,}25 à l'époque~90,
après être descendue à 1{,}83), alors que l'exécution corrigée décroît puis
oscille dans une bande étroite. Un écart de niveau seul serait compatible avec
un simple manque de capacité; la différence de régularité indique que ce n'est
pas l'optimisation qui échoue, mais la cible qui est par endroits incohérente.

\begin{figure}[htbp]
\centering
\begin{tikzpicture}
\begin{axis}[
    width=0.86\textwidth, height=6.2cm,
    xlabel={époque},
    ylabel={$\mathcal{L}_{\mathrm{val}}$},
    xmin=0, xmax=205, ymin=0, ymax=4.6,
    grid=both, grid style={black!12},
    legend pos=north east, legend cell align=left,
    legend style={font=\small, fill=white, draw=black!30},
    tick label style={font=\small}, label style={font=\small}
]
\addplot[thick, black!55, mark=none] coordinates {
(1,4.4132) (5,3.1977) (10,2.9846) (15,2.8622) (20,1.9178) (25,2.5414)
(30,1.8580) (35,1.8264) (40,1.8411) (45,1.6704) (50,2.4013) (55,1.8430)
(60,1.8795) (65,1.8297) (70,2.0997) (75,1.7269) (80,1.6235) (85,1.6767)
(90,2.2544) (95,1.7458) (100,1.8017) (105,1.5887) (110,1.6473) (115,1.5979)
(120,1.6500) (125,1.5593) (130,1.7071) (135,1.7752) (140,1.6060) (145,1.5262)
(150,1.7044) (155,1.5411) (160,1.6706) (165,1.7199) (170,1.5688) (175,1.5578)
(180,1.6930) (185,1.4903) (190,1.6666) (195,1.8044) (200,1.5840)
};
\addlegendentry{\code{isaac-first-real} (avant filtrage)}
\addplot[very thick, black, mark=none] coordinates {
(1,0.8613) (5,0.6507) (10,0.4823) (15,0.4528) (20,0.5327) (25,0.4174)
(30,0.4253) (35,0.4477) (40,0.4028) (45,0.4336) (50,0.4491) (55,0.4278)
(60,0.4306) (65,0.4467) (70,0.4021) (75,0.4183) (80,0.4784) (85,0.4335)
(90,0.4310) (95,0.4498) (100,0.4224) (105,0.4161) (110,0.4300) (115,0.4219)
(120,0.4287) (125,0.4009) (130,0.4767) (135,0.4439) (140,0.4316) (145,0.4222)
(150,0.4297) (155,0.4353) (160,0.4222) (165,0.4462) (170,0.4250) (175,0.4021)
(180,0.4293) (185,0.4248) (190,0.4029) (195,0.4234) (200,0.4144)
};
\addlegendentry{\code{isaac-satfix} (après filtrage)}
\end{axis}
\end{tikzpicture}
\caption{Perte de validation~\eqref{eq:Ltotal} au fil des époques, relevée tous
les 5~pas, pour les deux exécutions sur données physiques. Le filtrage de
saturation~\eqref{eq:satfilter} abaisse le plateau d'un facteur~4 environ et
supprime les sursauts de la courbe initiale.}
\label{fig:courbes}
\end{figure}

La performance par axe est plus informative que la perte globale, car elle
s'exprime en newton-mètres et se compare aux limites de couple. La
figure~\ref{fig:rmse} donne l'erreur quadratique moyenne par axe pour les deux
exécutions, et le tableau~\ref{tab:perjoint} en tire les conséquences pour la
commande.

\begin{figure}[htbp]
\centering
\begin{tikzpicture}
\begin{axis}[
    width=0.86\textwidth, height=5.6cm,
    ybar, bar width=7pt,
    xlabel={axe articulaire},
    ylabel={RMSE de validation [Nm]},
    symbolic x coords={J1,J2,J3,J4,J5,J6,J7},
    xtick=data, ymin=0, ymax=2.0,
    grid=major, grid style={black!12},
    legend pos=north east, legend cell align=left,
    legend style={font=\small, fill=white, draw=black!30},
    nodes near coords, every node near coord/.append style={font=\tiny,
        /pgf/number format/fixed, /pgf/number format/precision=2,
        /pgf/number format/fixed zerofill},
    tick label style={font=\small}, label style={font=\small}
]
\addplot[fill=black!25, draw=black!50] coordinates {
(J1,1.5694) (J2,1.4355) (J3,1.1351) (J4,1.7685) (J5,0.9680) (J6,0.4885)
(J7,0.5180)};
\addlegendentry{\code{isaac-first-real}}
\addplot[fill=black!70, draw=black] coordinates {
(J1,0.7996) (J2,0.8823) (J3,0.8979) (J4,0.5396) (J5,0.3732) (J6,0.3593)
(J7,0.2986)};
\addlegendentry{\code{isaac-satfix}}
\end{axis}
\end{tikzpicture}
\caption{Erreur quadratique moyenne de prédiction du couple, par axe, sur la
partition de validation (14\,830 échantillons), avant et après le filtrage de
saturation~\eqref{eq:satfilter}. L'amélioration est présente sur les sept axes,
la plus forte sur l'axe~4 (facteur 3{,}3).}
\label{fig:rmse}
\end{figure}

\begin{table}[htbp]
\centering
\caption{Erreur de prédiction par axe pour l'exécution de référence
\code{isaac-satfix}, borne d'erreur retenue $\epsilon_j$ (quantile 99{,}9 de
$|\taupred - \taureal|$ sur la validation), et gains PD résultants
de~\eqref{eq:gains} avec $\kappa = 2$. La dernière ligne rapporte l'erreur
relative à la limite de couple de l'axe, seule grandeur comparable d'un axe à
l'autre (section~\ref{sec:res-quoi}).}
\label{tab:perjoint}
\small
\begin{tabular}{@{}lccccccc@{}}
\toprule
Axe $j$ & 1 & 2 & 3 & 4 & 5 & 6 & 7 \\
\midrule
RMSE [Nm] & 0{,}800 & 0{,}882 & 0{,}898 & 0{,}540 & 0{,}373 & 0{,}359
& 0{,}299 \\
$\epsilon_j$ [Nm] & 5{,}20 & 5{,}66 & 3{,}06 & 3{,}80 & 2{,}10 & 2{,}38
& 1{,}57 \\
$K_{d,jj}$ [Nm$\cdot$s$\cdot$rad$^{-1}$] & 10{,}40 & 11{,}32 & 6{,}12 & 7{,}60
& 4{,}20 & 4{,}76 & 3{,}14 \\
$K_{p,jj}$ [Nm$\cdot$rad$^{-1}$] & 27{,}04 & 32{,}04 & 9{,}36 & 14{,}44
& 4{,}41 & 5{,}66 & 2{,}46 \\
\midrule
RMSE $/\ \bar{\tau}_j$ [\%] & 0{,}92 & 1{,}01 & 1{,}03 & 0{,}62 & 3{,}11
& 2{,}99 & 2{,}49 \\
\bottomrule
\end{tabular}
\end{table}

L'erreur absolue est plus faible sur les trois derniers axes (0{,}30 à
0{,}37~Nm) que sur les quatre premiers (0{,}54 à 0{,}90~Nm), ce qui est attendu
puisque ces axes portent moins de charge et de couplage inertiel. Rapportée à la
limite de couple, la hiérarchie s'inverse: l'erreur relative est de 0{,}62 à
1{,}03\% sur les axes~1 à~4 mais de 2{,}49 à 3{,}11\% sur les axes~5 à~7. Ce
renversement n'est pas un détail de présentation: il détermine quel axe
consomme le plus d'autorité d'actionneur pour compenser l'erreur de modèle, et
il sera invoqué en section~\ref{sec:res-desequilibre} pour justifier de ne pas
corriger un déséquilibre mesuré.

\subsection{Les garanties physiques survivent-elles à l'entraînement~?}
\label{sec:res-physique}

Les contraintes ne sont pas seulement déclarées: elles sont mesurées à chaque
époque, et cette mesure sépare nettement les degrés de garantie introduits en
section~\ref{sec:residu}.

\parag{Degré~1, limites de couple.} Sur l'exécution de référence, la
violation maximale de la contrainte~\eqref{eq:gtorque} est de 0{,}0000~Nm à
chacune des 200~époques enregistrées: le modèle n'a jamais prédit un couple
hors limites sur les données d'entraînement. À titre de comparaison,
l'exécution avant filtrage présentait des violations intermittentes atteignant
22{,}75~Nm à l'époque~200, 4{,}56~Nm à l'époque~80 et 4{,}26~Nm à l'époque~110.
Ce n'est pas un défaut du mécanisme dual: le modèle prédisait des couples
irréalisables parce que les échantillons saturés lui présentaient des cibles
résiduelles absurdes.

\parag{Degré~1, dissipativité.} La violation moyenne de la
contrainte~\eqref{eq:gdissip} décroît de 0{,}1070 à l'époque~1 jusqu'à 0{,}0006
à l'époque~200 sur l'exécution de référence, avec une décroissance quasi
monotone (figure~\ref{fig:dissip}). Sur l'exécution avant filtrage, la même
quantité oscille sans se stabiliser et remonte jusqu'à 0{,}0914 à l'époque~110
et 0{,}0657 à l'époque~200. C'est le point d'interprétation le plus important
de cette sous-section: la contamination des données ne dégradait pas seulement
la précision, elle déstabilisait le mécanisme de montée duale lui-même. Une
garantie de degré~1 n'est donc pas seulement statistique, elle est aussi
\emph{fragile} vis-à-vis de la qualité des données.

\begin{figure}[htbp]
\centering
\begin{tikzpicture}
\begin{semilogyaxis}[
    width=0.86\textwidth, height=5.8cm,
    xlabel={époque},
    ylabel={violation moyenne de dissipativité},
    xmin=0, xmax=205, ymin=4e-4, ymax=0.3,
    grid=both, grid style={black!12},
    legend pos=south west, legend cell align=left,
    legend style={font=\small, fill=white, draw=black!30},
    tick label style={font=\small}, label style={font=\small}
]
\addplot[thick, black!55, mark=none] coordinates {
(1,0.1480) (5,0.0537) (10,0.0141) (15,0.0180) (20,0.0092) (25,0.0138)
(30,0.0117) (35,0.0610) (40,0.0100) (45,0.0049) (50,0.0082) (55,0.0133)
(60,0.0085) (65,0.0070) (70,0.0061) (75,0.0033) (80,0.0556) (85,0.0033)
(90,0.0214) (95,0.0051) (100,0.0028) (105,0.0053) (110,0.0914) (115,0.0043)
(120,0.0304) (125,0.0025) (130,0.0030) (135,0.0309) (140,0.0020) (145,0.0015)
(150,0.0029) (155,0.0029) (160,0.0080) (165,0.0014) (170,0.0029) (175,0.0012)
(180,0.0025) (185,0.0009) (190,0.0017) (195,0.0015) (200,0.0657)
};
\addlegendentry{\code{isaac-first-real} (avant filtrage)}
\addplot[very thick, black, mark=none] coordinates {
(1,0.1070) (5,0.0522) (10,0.0259) (15,0.0123) (20,0.0186) (25,0.0040)
(30,0.0083) (35,0.0105) (40,0.0081) (45,0.0041) (50,0.0021) (55,0.0055)
(60,0.0030) (65,0.0082) (70,0.0039) (75,0.0009) (80,0.0037) (85,0.0040)
(90,0.0066) (95,0.0036) (100,0.0068) (105,0.0052) (110,0.0060) (115,0.0011)
(120,0.0018) (125,0.0033) (130,0.0033) (135,0.0035) (140,0.0023) (145,0.0053)
(150,0.0023) (155,0.0012) (160,0.0068) (165,0.0006) (170,0.0099) (175,0.0026)
(180,0.0012) (185,0.0029) (190,0.0013) (195,0.0017) (200,0.0006)
};
\addlegendentry{\code{isaac-satfix} (après filtrage)}
\end{semilogyaxis}
\end{tikzpicture}
\caption{Violation moyenne de la contrainte de dissipativité~\eqref{eq:gdissip}
au fil des époques, échelle logarithmique. Après filtrage des échantillons
saturés, la montée duale converge vers $6\cdot 10^{-4}$; avant filtrage, elle
oscille et remonte jusqu'à $9{,}1\cdot 10^{-2}$.}
\label{fig:dissip}
\end{figure}

\parag{Degré~2, garantie structurelle.} Indépendamment de ce qui précède, la
propriété~\eqref{eq:dissip} du module de frottement est vérifiée par
construction et testée numériquement sur des états tirés aléatoirement, y
compris hors de la distribution d'entraînement. Les coefficients de dissipation
appris sont strictement positifs, avec des valeurs moyennes par axe de l'ordre
de 1{,}17 à 1{,}63 sur la campagne multi-charge. C'est le seul résultat du
travail qui soit vrai pour tout état et tout jeu de poids, et c'est ce qui
justifie l'existence du degré~2 malgré le bon comportement du degré~1 sur
données saines.

\parag{Un effet secondaire attendu.} Lorsque le module de frottement est
actif, le multiplicateur $\lambda_d$ croît plus lentement, puisqu'une partie de
la contrainte de dissipativité est absorbée structurellement plutôt que
pénalisée. Ce comportement est intentionnel et cohérent, mais il implique
qu'une comparaison de $\lambda_d$ entre exécutions avec et sans le module n'est
pas interprétable telle quelle — remarque à retenir pour toute ablation
ultérieure.

\subsection{Une hypothèse testée de façon contrôlée: la contamination par
saturation}
\label{sec:res-satfix}

C'est le seul résultat de ce travail obtenu par une expérience contrôlée avec
hypothèse préalable, et il mérite d'être détaillé pour cette raison.

\parag{Observation initiale.} L'outil de mesure de borne d'erreur, écrit pour
alimenter~\eqref{eq:gains}, a révélé sur la première exécution une erreur de
prédiction \emph{maximale} par axe de 91 à 110~Nm — supérieure à la limite de
couple de 87~Nm de l'axe lui-même. Une erreur de modèle plus grande que le
couple physiquement réalisable n'est pas un défaut de modèle: c'est le signe
que la cible d'apprentissage est fausse pour certains échantillons. C'est
exactement la situation contre laquelle Deng \emph{et al.}~\cite{deng2024}
avertissent, et elle serait restée invisible sur la seule erreur moyenne.

\parag{Hypothèse.} L'inspection des pires échantillons a montré que leur
couple mesuré se situait presque exactement sur la limite de couple.
L'hypothèse formulée était donc: l'actionneur simulé sature à son effort
maximal sous un suivi de trajectoire agressif; $\taureal$ vaut alors la valeur
d'écrêtage, alors que $\taurbd$ est calculé depuis l'accélération de référence,
en supposant la trajectoire atteinte. Le résidu cible de ces échantillons est
donc un artefact.

\parag{Test et résultat.} Le critère~\eqref{eq:satfilter} a été appliqué, les
trois jeux de données régénérés et l'entraînement relancé à réglages
strictement identiques. Le tableau~\ref{tab:exp} résume l'effet mesuré. Le
rapport entre le gain et le coût est ce qui confirme l'hypothèse: 456
échantillons écartés, soit 0{,}31\% du jeu de données, valent un facteur 3{,}7
sur la perte de validation, un facteur~110 sur la violation de dissipativité
finale, et l'annulation complète des violations de limite de couple. Un très
petit nombre d'échantillons corrompus dominait donc l'optimisation. C'est
également le seul endroit du travail où une variable a été modifiée seule, tout
le reste étant tenu constant, ce qui autorise une conclusion causale.

\begin{table}[htbp]
\centering
\caption{Expérience contrôlée sur le filtrage de saturation~\eqref{eq:satfilter}:
deux entraînements identiques en tous points sauf le filtre appliqué aux
données. $\mathcal{L}_{\mathrm{val}}$ est la perte totale~\eqref{eq:Ltotal}
évaluée sur la partition de validation; les violations sont celles
de~\eqref{eq:gtorque} et~\eqref{eq:gdissip} relevées à chaque époque.}
\label{tab:exp}
\small
\begin{tabular}{@{}lrrr@{}}
\toprule
Grandeur mesurée & Avant filtrage & Après filtrage & Variation \\
\midrule
Identifiant d'exécution & \code{isaac-first-real} & \code{isaac-satfix} & — \\
Échantillons retenus & 148\,760 & 148\,304 & $-456$ ($-0{,}31$\%) \\
Époques d'entraînement & 200 & 200 & inchangé \\
Meilleure $\mathcal{L}_{\mathrm{val}}$ & 1{,}4767 & 0{,}3995
& $\div\,3{,}70$ \\
Époque du meilleur point & 177 & 144 & — \\
RMSE moyenne axes 1--4 [Nm] & 1{,}477 & 0{,}780 & $-47$\% \\
RMSE moyenne axes 5--7 [Nm] & 0{,}658 & 0{,}344 & $-48$\% \\
Violation de dissipativité, dernière époque & 0{,}0657 & 0{,}0006
& $\div\,110$ \\
Violation de couple maximale observée [Nm] & 22{,}75 & 0{,}00 & annulée \\
\bottomrule
\end{tabular}
\end{table}

\subsection{De l'erreur de modèle aux gains: quelle borne, quel certificat~?}
\label{sec:res-bound}

Les gains de commande dépendent du choix de la borne $\epsilon_j$
dans~\eqref{eq:gains}, et ce choix mérite d'être exposé plutôt que subi. Après
filtrage, la queue de distribution de l'erreur de prédiction reste
asymétrique: le quantile 99{,}9 est de 3 à 6~Nm sur les axes~1 à~4, mais le
maximum absolu atteint 26 à 65~Nm sur environ 0{,}09\% des échantillons de
validation. L'origine probable de ces valeurs extrêmes résiduelles est un
transitoire de correction du servo simulé en début de segment de trajectoire, et
non un défaut du modèle appris.

Le choix retenu est de prendre pour $\epsilon_j$ le quantile 99{,}9 et non le
maximum. C'est un compromis assumé, et il doit être présenté comme tel:
injecter le maximum brut dans~\eqref{eq:gains} produirait des gains de l'ordre
de $K_d \approx 130$ et $K_p \approx 4000$, incompatibles avec l'exigence de
couple lisse posée en section~\ref{sec:degre0}, pour un bénéfice négligeable
puisque l'écrêtage aux limites d'actionneur constitue déjà une protection
indépendante contre ces rares transitoires. La conséquence logique est un
affaiblissement du certificat: la garantie obtenue n'est pas \og valable pour
tout échantillon \fg{} mais \og valable pour 99{,}9\% des conditions de
fonctionnement observées, avec les limites matérielles comme filet pour le
reste \fg{}. Les gains résultants figurent au tableau~\ref{tab:perjoint}.

\subsection{Le modèle appris pilote-t-il effectivement le robot~?}
\label{sec:res-loop}

\parag{Planification et exécution.} La planification MoveIt2 et l'exécution en
mode position ont été validées: le bras simulé sous MuJoCo suit les plans
demandés, et les vues MuJoCo et RViz restent cohérentes entre elles. Ce jalon
est stable et reproductible, mais il ne démontre rien sur le modèle appris,
puisque c'est l'exécuteur en position de MoveIt2 qui commande.

\parag{Suivi de trajectoire par le modèle appris.} Le jalon décisif est donc le
routage de la trajectoire planifiée vers le n\oe ud de commande, de sorte que ce
soit le couple issu de~\eqref{eq:ctpd} — donc le modèle appris — qui pilote le
robot. Ce jalon a été atteint et confirmé numériquement: sur une fenêtre de
8~s bracketant un mouvement, chacun des sept axes s'est déplacé de 0{,}2 à
1{,}15~rad en convergeant vers sa cible. Ce résultat est qualitativement
différent du défaut qu'il a fallu lever pour l'obtenir, où la position mesurée
était figée au bit près (minimum égal au maximum en précision flottante double)
sur toute la fenêtre de capture, alors que le couple commandé était correct et
variait dans le temps — signature qui reviendra en
section~\ref{sec:disc-freeze}.

\parag{Décomposition du couple commandé.} Le tableau~\ref{tab:tau} donne une
capture réelle de la décomposition de~\eqref{eq:ctpd} pendant une phase de
suivi, relevée à un pas de commande sur trois axes instrumentés. Ces valeurs
documentent la hiérarchie effective des termes et valident la décomposition
choisie en~\eqref{eq:decompo}: sur l'axe~4, le terme calculé domine largement
(22{,}4~Nm, essentiellement de la gravité dans cette configuration), le résidu
appris est faible (0{,}03~Nm) et la correction PD est intermédiaire
($-4{,}27$~Nm); sur l'axe~7, dont la limite est de 12~Nm, le terme calculé est
nul, tandis que le résidu appris ($-0{,}10$~Nm) et la correction PD (0{,}78~Nm)
sont du même ordre de grandeur. Autrement dit, la part relative du terme appris
croît là où le terme rigide s'efface, ce qui est le comportement attendu d'une
architecture grise.

\begin{table}[htbp]
\centering
\caption{Décomposition du couple commandé~\eqref{eq:ctpd} relevée en
fonctionnement, à un pas de commande donné, sur trois axes instrumentés. La
colonne \og écrêté \fg{} montre que l'écrêtage aux limites n'était pas actif à
cet instant.}
\label{tab:tau}
\small
\begin{tabular}{@{}lrrrrr@{}}
\toprule
Axe & $\taurbd$ [Nm] & $\taures$ [Nm] & $K_p e + K_d\dot{e}$ [Nm]
& somme [Nm] & écrêté [Nm] \\
\midrule
4 & $+22{,}352$ & $+0{,}026$ & $-4{,}266$ & $+18{,}112$ & $+18{,}112$ \\
6 & $+2{,}085$  & $+0{,}070$ & $+1{,}682$ & $+3{,}837$  & $+3{,}837$ \\
7 & $-0{,}000$  & $-0{,}095$ & $+0{,}782$ & $+0{,}687$  & $+0{,}687$ \\
\bottomrule
\end{tabular}
\end{table}

\parag{Erreur de suivi résiduelle.} La loi~\eqref{eq:ctpd} est purement
proportionnelle-dérivée, sans terme intégral. Une erreur statique subsiste donc
sous charge de gravité, et elle a été mesurée: sur un mouvement vers une pose
de référence, six axes sur sept convergent à moins de 0{,}04~rad de leur cible,
tandis que l'axe~3 se stabilise sur un résidu d'environ 0{,}14~rad. Ce
comportement est conforme à ce qu'on attend d'un correcteur PD sans intégrateur
et n'est pas un défaut de mise en \oe uvre; il a conduit à relâcher le seuil de
convergence de 0{,}05 à 0{,}15~rad sur la base de cette mesure, valeur qui reste
bien inférieure aux 0{,}3~rad observés lors d'un échec authentique et ne masque
donc pas les vraies défaillances.

\parag{Cadence.} Le n\oe ud est configuré et exécuté à
\code{control_rate_hz = 1000}, ce qui est confirmé par sa trace de démarrage. La
cadence nominale est donc atteinte. Ce qui n'est \emph{pas} établi est la marge
temporelle réelle: aucune mesure de latence d'inférence n'a été réalisée, donc
rien ne prouve pour l'instant que le budget de 1~ms soit respecté avec marge
plutôt que tout juste. C'est la principale lacune de validation de la
contribution~2, et la décomposition de latence publiée par
Agrawal \emph{et al.}~\cite{agrawal2026} indique la forme exacte que devrait
prendre la mesure manquante (section~\ref{sec:futur}).

\subsection{Le test d'intégration: la tâche de préhension}
\label{sec:res-task}

Les briques individuelles sont validées séparément: la pince possède une
physique réelle (glissières prismatiques, contraintes d'égalité de
synchronisation, coussinets de contact aux extrémités) et exécute des commandes
d'ouverture et de fermeture avec un résultat d'action confirmé; la scène et
l'objet sont correctement rendus et positionnés; le client de mouvement
cartésien planifie et exécute des poses cibles.

En revanche, \textbf{la séquence de préhension complète n'a jamais abouti}.
Douze exécutions instrumentées ont été enregistrées avec journalisation
complète; les huit dernières, correspondant à l'état le plus avancé du code,
sont résumées au tableau~\ref{tab:stage4}. Toutes se terminent par le même
verdict — dépassement du délai de convergence du bras lors d'un sous-mouvement —
avant même que la pince n'atteigne l'objet.

\begin{table}[htbp]
\centering
\caption{Huit exécutions instrumentées de la séquence de préhension. L'erreur
finale est la distance cartésienne de la bride à sa cible planifiée, mesurée par
cinématique directe. Toutes les exécutions échouent sur le même verdict
(dépassement du délai de convergence), mais avec des amplitudes d'erreur et des
axes fautifs différents; c'est cette dispersion qui est diagnostique.}
\label{tab:stage4}
\small
\begin{tabular}{@{}lccclc@{}}
\toprule
Exécution & Erreur finale [m] & Tolérance [m] & Rapport & Axe le plus fautif
& Verdict \\
\midrule
2  & 0{,}2464 & 0{,}030 & 8{,}2 & \code{panda_joint7} & délai dépassé \\
9  & 0{,}0824 & 0{,}030 & 2{,}7 & \code{panda_joint4} & délai dépassé \\
10 & 0{,}0846 & 0{,}030 & 2{,}8 & \code{panda_joint7} & délai dépassé \\
11 & 0{,}0858 & 0{,}030 & 2{,}9 & \code{panda_joint4} & délai dépassé \\
12 & 0{,}0322 & 0{,}030 & 1{,}1 & \code{panda_joint5} & délai dépassé \\
13 & 0{,}0991 & 0{,}030 & 3{,}3 & \code{panda_joint4} & délai dépassé \\
14 & 0{,}0822 & 0{,}040 & 2{,}1 & \code{panda_joint4} & délai dépassé \\
15 & 0{,}1009 & 0{,}040 & 2{,}5 & \code{panda_joint7} & délai dépassé \\
\bottomrule
\end{tabular}
\end{table}

Le résultat le plus important de cette série n'est pas le taux d'échec de
100\%, mais la \emph{dispersion} des échecs, qui distingue deux phénomènes de
nature différente. L'exécution~12 échoue à 0{,}0322~m pour une tolérance de
0{,}0300~m, c'est-à-dire à 7\% près sur l'axe~5: c'est un échec de précision
statique de suivi, entièrement cohérent avec l'erreur statique du correcteur PD
documentée en section~\ref{sec:res-loop}, et sans mystère. Les exécutions 9
à~11, 13 à~15 échouent au contraire à 0{,}08--0{,}10~m, c'est-à-dire d'un
facteur 2 à 3 au-delà de la tolérance, et — c'est le point décisif — avec un axe
fautif dont la position mesurée n'évolue \emph{plus du tout} pendant les huit
secondes d'attente de convergence. La trace de l'exécution~15 est sans
ambiguïté: \code{panda_joint7} est relevé à $-0{,}7853$~rad pour une cible de
$-0{,}4428$~rad, soit une erreur de $-0{,}3425$~rad, et cette valeur est
identique aux quatre décimales à chacun des huit relevés successifs espacés
d'une seconde, pendant que l'erreur cartésienne stagne à 0{,}1009~m. Il ne
s'agit donc pas d'une convergence lente: l'axe est immobile.

Cette distinction entre \og échec de précision \fg{} et \og axe immobile \fg{}
oriente tout le diagnostic. Le premier phénomène est expliqué et borné; le
second n'est pas résolu, et son investigation constitue le principal résultat
négatif de ce travail (section~\ref{sec:disc-freeze}).

% =====================================================================
\section{Discussion}
\label{sec:discussion}
% =====================================================================

\subsection{Ce qui se comporte comme prévu}

Trois éléments se comportent conformément à la conception, et il vaut la peine
de dire en quoi précisément. D'abord l'architecture grise: une erreur de
prédiction de 0{,}30 à 0{,}90~Nm par axe sur des données de dynamique d'un
moteur physique, obtenue avec un réseau modeste, confirme que déléguer la
dynamique du corps rigide à RNEA laisse au réseau une tâche bien conditionnée —
ce que corrobore la décomposition du tableau~\ref{tab:tau}, où le résidu ne
prend de l'importance relative que là où le terme rigide s'annule. Ensuite les
garanties physiques, aux trois degrés: violation de limites de couple nulle sur
les 200~époques, violation de dissipativité ramenée à $6\cdot 10^{-4}$, et
garantie structurelle inconditionnelle sur le terme de frottement. Enfin la
double prise en compte de la charge, qui n'a exigé aucun ajustement particulier
pour couvrir \SIlist{0;1;3}{kg} dans un modèle unique, ce qui est l'indice le plus
direct que la frontière tracée en~\eqref{eq:decompo} — charge à l'analytique,
état de mouvement au résidu — était la bonne.

\subsection{Deux défauts de données ont pesé plus que tout réglage}
\label{sec:disc-urdf}

Deux défauts de qualité de données ont eu, à eux seuls, plus d'effet sur les
résultats que tout réglage d'hyperparamètre effectué dans ce travail. Le premier
est l'absence de balises inertielles dans l'URDF de référence: RNEA a
fonctionné pendant toute la première partie du projet sur des inerties de
remplacement de la bibliothèque, et non sur les propriétés massiques réelles du
Franka. Ce défaut n'a été découvert que par comparaison avec un moteur physique
disposant de sa propre source de masses. Il a été corrigé à partir des valeurs
officielles de la description du constructeur, puis \emph{indépendamment
recoupé} par une troisième source: les masses et inerties du modèle MJCF
distribué par le dépôt de modèles MuJoCo, issues de la même description amont,
coïncident avec les valeurs corrigées. Ce recoupement à trois sources est le
seul moyen trouvé de se prémunir contre une erreur de ce type, qui ne produit
aucun symptôme visible et dégrade silencieusement le terme dominant du modèle.

Le second est le filtrage de saturation, dont l'effet est quantifié au
tableau~\ref{tab:exp}: 0{,}31\% d'échantillons corrompus coûtaient un
facteur~3{,}7 sur la perte de validation et déstabilisaient la montée duale.
L'enseignement méthodologique dépasse ce projet: dans une approche informée par
la physique, un échantillon dont la cible viole la physique n'est pas simplement
du bruit, il est activement nuisible, car les contraintes de la perte tentent de
s'y conformer. Une donnée aberrante coûte donc davantage à un modèle contraint
qu'à un modèle libre — conclusion qui va à l'encontre de l'intuition selon
laquelle les contraintes protègent du bruit.

\subsection{Un déséquilibre par axe réel mais laissé sans correction}
\label{sec:res-desequilibre}

L'asymétrie des limites de couple du tableau~\ref{tab:constantes} laisse
attendre un déséquilibre d'optimisation entre les axes~1 à~4 et les axes~5
à~7: le terme de fidélité~\eqref{eq:Ldata} étant une erreur quadratique non
pondérée en newton-mètres, il accorde mécaniquement plus de poids aux axes
capables de couples élevés. Un diagnostic de suivi a donc été implémenté, qui
rapporte le rapport entre l'erreur moyenne des trois derniers axes et celle des
quatre premiers à la fin de chaque campagne.

Le résultat mesuré est instructif à double titre. Sur toutes les campagnes à
résidu synthétique, ce rapport valait 1{,}0 — aucun déséquilibre. Sur les deux
campagnes sur données physiques, il vaut 0{,}4 (0{,}780~Nm de moyenne sur les
axes~1 à~4 contre 0{,}344~Nm sur les axes~5 à~7), et cette valeur est
\emph{identique avant et après} le filtrage de saturation. Le déséquilibre est
donc, premièrement, un signal réel de la dynamique du moteur physique et non un
artefact de données contaminées; et, deuxièmement, un signal invisible sur
données synthétiques, ce qui constitue une raison supplémentaire de ne pas
valider une architecture sur résidu inventé. Il reste néanmoins en dessous du
seuil qui déclencherait une repondération par moyenne mobile exponentielle des
résidus, et aucune correction n'a été appliquée. La lecture en erreur relative
du tableau~\ref{tab:perjoint} montre du reste que le déséquilibre s'inverse
lorsqu'on normalise par la limite de couple; c'est un argument de plus pour ne
pas corriger \emph{à l'aveugle}, sans un critère explicite portant sur la
grandeur qui compte réellement pour la commande.

Il faut cependant se garder de présenter cette décision comme mieux établie
qu'elle ne l'est. Le diagnostic ci-dessus \emph{constate} le déséquilibre a
posteriori~; il ne \emph{teste} pas la correction qu'il rend pourtant
évidente~: aucune exécution avec une matrice de pondération diagonale
$\mathrm{diag}(1/\bar{\tau}_j^2)$ dans~\eqref{eq:Ldata} — la repondération
directe suggérée par le seuil relatif du tableau~\ref{tab:perjoint} — n'a été
menée à réglages par ailleurs identiques. L'argument \og le rapport
s'inverse en relatif \fg{} justifie de ne pas corriger \emph{sans réfléchir},
mais ne démontre pas que ne pas corriger soit la meilleure décision~: il se
peut qu'une perte pondérée réduise l'erreur relative sur les axes~5 à~7 sans
dégrader sensiblement celle des axes~1 à~4, auquel cas l'absence de correction
serait un choix sous-optimal plutôt que neutre. Cette lacune de rigueur est
reconnue explicitement ici et reportée comme travail futur bien délimité
(section~\ref{sec:futur})~: exécuter la même campagne avec perte pondérée par
$1/\bar{\tau}_j^2$, toutes choses égales par ailleurs, et comparer l'erreur
relative résultante à celle du tableau~\ref{tab:perjoint}.

\subsection{Le résultat négatif principal: l'immobilisation intermittente
d'un axe}
\label{sec:disc-freeze}

Le phénomène décrit en section~\ref{sec:res-task} — un ou plusieurs axes dont la
position mesurée cesse d'évoluer alors que le couple commandé reste correct —
n'est pas résolu. L'investigation est présentée ici parce que son étendue et sa
méthode ont une valeur en elles-mêmes, et parce que les hypothèses écartées le
sont sur des bases empiriques solides et non par élimination faute de mieux.

Ont été éliminées expérimentalement: la collision avec l'objet saisi (le cube a
été déplacé à 2{,}5~m hors de portée: le symptôme persiste à l'identique); un
défaut d'ordre ou de nommage des articulations (vérifié indépendamment en trois
endroits du code, et l'URDF déclare bien les axes dans l'ordre canonique); un
écrêtage silencieux aux limites d'actionneur (les couples relevés restent
largement dans la plage, cf. tableau~\ref{tab:tau}); un verrouillage du mode de
commande des actionneurs par l'intergiciel (le journal du n\oe ud de commande
montre les sept axes en mode couple, vérifié trois fois, et la lecture du code
source de la version exactement installée confirme une liaison par nom, uniforme
sur les sept axes); et enfin la loi de commande elle-même. Cette dernière
élimination est la plus concluante: l'axe fautif présente une erreur de
position réelle et persistante avec une correction PD de signe correct et
d'amplitude substantielle ($-4$ à $-11$~Nm selon la cible), et sous \emph{trois}
régimes de couple différents — couple calculé nominal d'environ 18~Nm, une
exécution à couple nul confirmée, et un couple constant de 40~Nm publié
manuellement en contournant tout le code de commande — la position mesurée varie
de moins de $10^{-4}$~rad. Le phénomène se situe donc sous la couche
ROS2/Python, et aucune modification de la loi de commande ou du modèle appris ne
peut l'atteindre.

Une cause réelle a été identifiée et corrigée: MuJoCo calcule les collisions
sur des enveloppes convexes, plus grosses que les maillages fins à partir
desquels la liste d'exclusions d'auto-collision avait été dérivée; certaines
paires (main et doigts contre l'avant-bras) pouvaient donc s'interpénétrer dans
des poses de poignet replié, produisant un contact parasite suffisamment raide
pour absorber 40~Nm et bloquer précisément les axes qui rapprochent la main de
l'avant-bras. La désactivation de l'auto-collision robot-robot, en laissant
intactes les collisions avec le sol et l'objet, a produit deux séquences
complètes de sous-mouvements convergeant normalement. \textbf{Mais le symptôme
est réapparu lors des tentatives suivantes avec le correctif actif}, ce qui
établit que cette cause, bien que réelle, n'explique pas toutes les occurrences.
L'enregistrement du sol comme obstacle de planification, tenté comme seconde
atténuation indépendante, n'a pas non plus empêché la récurrence
(exécution~15). Il faut donc rapporter ce correctif pour ce qu'il est: une
cause vérifiée mais insuffisante, et non une solution.

L'hypothèse restante à l'issue de ce qui précède était que le solveur de
cinématique inverse de MoveIt2 (OMPL/KDL) retourne des solutions articulaires
différentes pour une même cible cartésienne d'un appel à l'autre, et
qu'aucune des corrections précédentes ne contraint ce choix. Il faut la
formuler avec précaution~: la planification produit une consigne
\emph{avant} l'exécution, elle ne peut donc pas être la cause directe d'un
blocage qui se manifeste \emph{pendant} le suivi, avec un couple de
correction actif et correctement signé, comme établi plus haut. Ce que cette
hypothèse peut en revanche expliquer, c'est que des appels successifs à
MoveGroup pour une même cible cartésienne n'aboutissent pas nécessairement à
la même consigne articulaire, de sorte que des tentatives réelles
successives peuvent converger tantôt vers une configuration bénigne, tantôt
vers une configuration où un phénomène d'exécution encore non identifié se
déclenche.

Ce point a été testé directement, après la clôture de la campagne
d'exécutions instrumentées ci-dessus, par un diagnostic dédié qui élimine la
variable de planification~: plutôt que de laisser MoveGroup planifier (et
donc invoquer OMPL/KDL, susceptible de retourner l'une de plusieurs solutions
valides pour une même cible), une unique configuration articulaire est
calculée une fois pour toutes par une résolution numérique déterministe de la
cinématique inverse (moindres carrés amortis sous Pinocchio, sur le même URDF
que le reste de la chaîne, sans composante aléatoire), puis rejouée telle
quelle cinq fois de suite, avec réinitialisation complète du monde entre
chaque tentative. Le correctif d'auto-collision de la section précédente est
resté actif et confirmé présent tout au long de cet essai. Résultat~: les
cinq tentatives se sont figées de façon identique sur les axes~4, 6 et~7,
avec des erreurs de position reproductibles au bit près d'une tentative à
l'autre — par exemple l'axe~6 relevé à exactement $-0{,}7966$~rad, sans
dérive mesurable, de $t=2{,}2$~s à $t=7{,}2$~s, dans chacune des cinq
tentatives.

Ce résultat précise le diagnostic qui précède sans le contredire. Retirer la
variable de choix de solution ne rend pas le phénomène aléatoire~: il le rend
parfaitement déterministe, pour une configuration articulaire donnée. Ceci
réconcilie l'intermittence observée au tableau~\ref{tab:stage4} avec
l'élimination de la loi de commande comme cause~: les tentatives réelles de
préhension ne rejouent pas la même consigne articulaire d'une exécution à
l'autre, puisque OMPL/KDL choisit à chaque appel une solution parmi
plusieurs pour la même cible cartésienne, et seule une partie de ces
solutions correspond à une configuration qui se fige de façon déterministe
une fois atteinte. L'intermittence vivait donc entièrement dans le choix de
la configuration retenue par la planification, et non dans une quelconque
variabilité de la couche d'exécution elle-même, qui se révèle au contraire
rigoureusement reproductible à consigne fixée. Il ne faut cependant pas
sur-interpréter ce résultat en réhabilitant le solveur de cinématique inverse
comme cause du blocage~: ce qu'il établit est que la planification détermine
\emph{si} une exécution atteint une configuration défaillante, non ce qui
rend cette configuration défaillante une fois atteinte — cette seconde
question reste entièrement du ressort de l'exécution, et son mécanisme précis
(un contact non couvert par la liste d'exclusions, une saturation
d'actionneur, un phénomène propre au solveur physique de MuJoCo) n'est pas
encore isolé. La visualisation des contacts pendant une reproduction sur
cette configuration désormais fixée et reproductible à volonté est l'étape
suivante, planifiée mais non encore réalisée (section~\ref{sec:futur}).

\subsection{Une leçon sur l'endroit où se trouvent les défauts}

Sur l'ensemble du travail d'intégration, une quinzaine de défauts réels ont été
trouvés et corrigés, et leur répartition est plus instructive que leur nombre:
aucun ne se trouvait dans les mathématiques de la commande ou de
l'apprentissage, et la quasi-totalité se situait aux interfaces — ordre des
articulations dans un message, péremption d'une consigne, cadence d'une boucle
héritée, enveloppes convexes d'un moteur de collision, image clé d'un modèle de
simulation réinitialisant silencieusement la position d'un objet. C'est
exactement le constat que dresse Ibarz \emph{et al.}~\cite{ibarz2021} pour
l'apprentissage sur robots réels, où les difficultés dominantes ne sont pas
celles que la littérature algorithmique anticipe. Deux règles générales en ont
été tirées, applicables au-delà de ce projet: ne jamais indexer un message
d'état articulaire autrement que par nom, et ne jamais proposer de correctif sur
la base d'une description verbale d'un symptôme sans avoir capturé la trace
brute correspondante — dans ce travail, un symptôme uniquement décrit oralement
s'est révélé, une fois journalisé, être un défaut entièrement différent et
beaucoup plus important.

\subsection{Limites}
\label{sec:limites}

Les limites suivantes sont énoncées sans atténuation, car elles bornent la
portée de tout ce qui précède.

\begin{description}
    \item[Aucune validation sur robot physique.] La totalité des résultats est
    obtenue en simulation, avec deux simulateurs distincts pour les données
    (Isaac~Sim) et pour la boucle de commande (MuJoCo), dont l'écart de
    domaine (\emph{sim-to-sim gap}, section~\ref{sec:exec}) reste, lui aussi,
    non quantifié~: la validation en boucle fermée de la
    section~\ref{sec:res-loop} démontre que le modèle appris pilote un
    simulateur physique différent de celui où il a été identifié, ce qui est
    un indice de robustesse mais ne mesure pas directement sa précision sous
    MuJoCo. La contribution~4 (transfert simulation-réalité) est implémentée
    et vérifiée sur le plan physique, mais elle n'est pas validée: aucun
    enregistrement de mouvements exploratoires sur le Franka réel n'a été
    réalisé, donc l'équation~\eqref{eq:finetune} n'a jamais été résolue sur
    des données réelles.
    \item[Couverture éparse de l'espace articulaire.] Les données
    d'entraînement (section~\ref{sec:donnees}) sont générées autour de dix
    configurations centrales seulement, tirées par une séquence de Sobol dans
    un espace à 7~dimensions. Même bien réparties, dix configurations laissent
    de larges régions de l'espace articulaire sans couverture directe~: le
    modèle appris n'est identifié, à proprement parler, que dans un voisinage
    de ces dix configurations et de l'excursion locale des trajectoires de
    Fourier qui les entourent, et non sur l'espace atteignable dans son
    ensemble. Aucune capacité de généralisation à des régions non couvertes
    n'est revendiquée dans ce rapport.
    \item[Latence d'inférence non mesurée.] La boucle est cadencée à \SI{1000}{Hz},
    mais le temps d'exécution d'un pas de~\eqref{eq:ctpd} n'a pas été
    instrumenté. La contribution~2 est donc étayée par un fonctionnement
    observé, non par une marge temporelle mesurée.
    \item[Comparaison externe non homogène.] Liu \emph{et al.}~\cite{liu2024}
    rapportent leur erreur en rad$^{2}$ sur un rollout de 2~s; ce travail
    rapporte une erreur de couple en Nm. Les deux grandeurs ne sont pas
    convertibles l'une dans l'autre, donc \emph{aucune} comparaison numérique
    directe avec la référence n'est revendiquée dans ce rapport. Implémenter
    leur métrique de rollout est une condition préalable à toute affirmation de
    supériorité chiffrée.
    \item[Portée restreinte aux trajectoires en distribution.] Suivant le
    résultat négatif de Prabhakar \emph{et al.}~\cite{prabhakar2026}, selon
    lequel les contraintes physiques dégradent l'extrapolation temporelle de 6
    à 10\%, aucune capacité d'extrapolation hors distribution n'est revendiquée.
    \item[Certificat de stabilité affaibli.] Comme expliqué en
    section~\ref{sec:res-bound}, la borne d'erreur employée est un quantile
    99{,}9 et non un maximum; la garantie n'est donc pas un certificat au sens
    du pire cas.
    \item[Tâche de préhension non fonctionnelle.] La séquence complète n'a
    jamais abouti sur aucune des douze exécutions instrumentées
    (tableau~\ref{tab:stage4}). Le n\oe ud d'orchestration ROS2 correspondant
    n'existe pas non plus: seuls des scripts de test existent,
    délibérément, tant qu'une préhension fiable n'est pas obtenue.
    \item[Matrice d'ablation incomplète.] Le plan d'expériences prévoit des
    ablations systématiques (structure, dissipation diagonale contre Cholesky
    complète, choix d'activation, mode de contrainte, efficacité en données via
    la troncature du jeu d'apprentissage). Seul le drapeau de troncature est
    implémenté; les autres axes ne le sont pas, et l'effort a été redirigé
    vers l'intégration temps réel. Deux réglages en particulier sont fixés par
    convention plutôt que par étude~: le poids de pénalisation $\rho = 1$ du
    lagrangien augmenté (section~\ref{sec:degre1}), dont aucune sensibilité
    n'a été mesurée alors que la littérature associe classiquement
    l'instabilité de la montée duale à son calibrage~; et l'absence de
    pondération diagonale $1/\bar{\tau}_j^2$ dans la perte de
    fidélité~\eqref{eq:Ldata}, dont le bénéfice sur le déséquilibre par axe
    mesuré en section~\ref{sec:res-desequilibre} n'a pas été testé.
    \item[Opacité d'une dépendance binaire.] Le choix du paquet
    d'interfaçage MuJoCo, motivé en section~\ref{sec:exec}, s'est retourné en
    obstacle de diagnostic: sa nature de binaire compilé a imposé de recourir
    au code source amont de la version exacte installée pour trancher des
    hypothèses qui, avec une dépendance sous forme de sources, auraient été
    vérifiables directement.
    \item[Sous-échantillonnage possible en régime d'ablation.] Pour de très
    petites tailles d'échantillon (moins d'environ 2\,000 tirés d'un jeu de
    25\,000), le tirage aléatoire peut sous-représenter les régions à forte
    vitesse; cette réserve doit accompagner toute étude d'efficacité en
    données.
\end{description}

% =====================================================================
\section{Conclusion et perspectives}
\label{sec:conclusion}
% =====================================================================

\subsection{Synthèse}

Ce travail a construit et instrumenté une chaîne allant du fichier de
description d'un manipulateur jusqu'à un robot effectivement commandé par un
modèle dynamique appris. Elle repose sur un choix structurant, formulé en
réponse à un verrou identifié dans la littérature: conserver intact le terme
analytique de la dynamique du corps rigide, évalué automatiquement depuis
l'URDF, et n'apprendre qu'un résidu contraint et conditionné par la charge
portée, en distinguant explicitement trois degrés de garantie — rien, imposé par
la perte, vrai par construction.

Sur le plan quantitatif, le modèle de référence atteint une erreur de prédiction
de couple par axe de 0{,}30 à 0{,}90~Nm sur une partition de validation de
14\,830 échantillons, soit 0{,}62 à 3{,}11\% de la limite de couple de chaque
axe, pour un modèle unique couvrant trois valeurs de charge (\SIlist{0;1;3}{kg}). Les
contraintes physiques sont satisfaites: violation de limite de couple nulle sur
les 200~époques, violation de dissipativité ramenée à $6\cdot 10^{-4}$, et
garantie structurelle de dissipativité vraie par construction sur le terme de
frottement. Ce modèle a été inséré dans une loi de couple calculé à \SI{1000}{Hz} dont
les gains sont dérivés d'un résultat de stabilité appliqué à une borne d'erreur
mesurée et non postulée, et le suivi de trajectoire piloté par le modèle appris
a été confirmé numériquement sur les sept axes.

Sur le plan méthodologique, deux résultats de qualité de données méritent d'être
retenus au-delà de ce projet: un fichier URDF dépourvu d'inerties réelles peut
fausser silencieusement un terme analytique pendant des mois sans produire aucun
symptôme, et 0{,}31\% d'échantillons contaminés par la saturation d'un
actionneur suffisent à dégrader la perte de validation d'un facteur~3{,}7 tout
en déstabilisant la convergence des contraintes — un modèle contraint est donc
plus sensible aux données aberrantes qu'un modèle libre, et non moins.

Enfin, ce rapport documente un résultat négatif: la tâche de préhension
complète n'aboutit pas, et l'immobilisation intermittente d'un axe qui l'en
empêche n'est qu'à moitié expliquée. Une cause réelle a été trouvée et corrigée,
sans suffire.

\subsection{Travaux futurs}
\label{sec:futur}

Les perspectives se hiérarchisent en fonction de ce qui bloque quoi.

\parag{Court terme — lever le verrou de la tâche complète.} Le test proposé en
section~\ref{sec:disc-freeze} — rejouer plusieurs fois de suite la
\emph{même} configuration articulaire défaillante en s'affranchissant du choix
de solution d'OMPL/KDL — a depuis été réalisé et confirme le déterminisme du
blocage à configuration fixée (cinq tentatives sur cinq, erreurs de position
reproductibles au bit près). La priorité immédiate se déplace donc vers deux
tâches distinctes. D'abord, isoler le mécanisme physique exact à la
configuration désormais reproductible à volonté — le correctif d'auto-collision
étant actif et confirmé insuffisant sur cet essai — par la visualisation des
contacts MuJoCo, technique qui s'est avérée la plus discriminante de tout le
travail de déverminage. Ensuite, et indépendamment du résultat de cette
visualisation, contraindre ou biaiser le choix de solution d'OMPL/KDL pour
éviter que la planification ne sélectionne cette configuration défaillante
devient un travail de planification bien délimité, à traiter en amont de la
commande, quel que soit le mécanisme physique sous-jacent une fois qu'il est
identifié.

\parag{Court terme — clore la validation de la contribution 2.}
Instrumenter le temps d'exécution d'un pas de~\eqref{eq:ctpd} et publier la
distribution de latence, avec ses quantiles, contre le budget de 1~ms. La forme
souhaitable de ce résultat est celle d'une décomposition par composant — appel
RNEA, inférence du résidu, interpolation de consigne, publication du message —
sur le modèle de ce que rapporte Agrawal \emph{et al.}~\cite{agrawal2026}, car
c'est le maillon lent qui détermine la faisabilité et non l'inférence seule.
C'est un travail de faible ampleur qui transformerait une affirmation en mesure.

\parag{Moyen terme — la validation externe.} Implémenter la métrique de
rollout de 2~s de Liu \emph{et al.}~\cite{liu2024} afin de rendre possible une
comparaison homogène avec la référence, et compléter la matrice d'ablation, en
priorité sur trois axes: dissipation diagonale contre factorisation complète,
pour quantifier ce que coûte réellement l'hypothèse de rang de
Wang \emph{et al.}~\cite{wang2025spel}~; efficacité en données, pour situer ce
travail par rapport au seuil de 25\,000 échantillons annoncé par la
référence~; et perte pondérée par $1/\bar{\tau}_j^2$ contre perte non pondérée
actuelle, pour trancher la question laissée ouverte en
section~\ref{sec:res-desequilibre} sur le déséquilibre par axe.
La méthode de Sutanto \emph{et al.}~\cite{sutanto2020} fournit le protocole
approprié pour la première: fixer une précision cible et mesurer le coût
d'apprentissage nécessaire pour l'atteindre, plutôt que comparer des précisions
finales à budget égal.

\parag{Moyen terme — fermer la boucle simulation-réalité.} Enregistrer un
jeu de mouvements exploratoires sur le Franka physique afin de résoudre
effectivement~\eqref{eq:finetune}. C'est la seule voie pour faire passer la
contribution~4 du statut d'implémentation vérifiée à celui de résultat validé.
Ce même jeu de données permettrait de rejouer le diagnostic de déséquilibre par
axe, qui n'a jamais atteint son seuil de déclenchement sur données simulées mais
pourrait le franchir sur données réelles
(section~\ref{sec:res-desequilibre}).

\parag{Plus long terme — deux extensions naturelles.} D'une part,
l'imposition semi-supervisée des contraintes sur des points de l'espace
articulaire non étiquetés, proposée par
Djeumou \emph{et al.}~\cite{djeumou2022}: elle étendrait la garantie de
dissipativité au-delà des trajectoires vues à l'entraînement, c'est-à-dire
qu'elle rapprocherait le degré~1 du degré~2 sans exiger de le rendre
structurel. D'autre part, l'identification en ligne de la charge portée, qui
rendrait $\delta$ observable en fonctionnement au lieu d'être fournie; la voie
par les courants moteurs de Li \emph{et al.}~\cite{li2025ral} n'est pas
transposable au Franka en raison de ses réducteurs, mais une estimation par
moindres carrés sur les couples articulaires mesurés l'est, et constitue le
préalable à un déploiement sur robot réel.

% =====================================================================
\printbibliography[title={Références}]

\end{document}
